\documentclass[../main.tex]{subfiles}
\begin{document}
%==========================================TIÊU ĐỀ TẬP========================================
\begin{table}[h]
  \begin{adjustwidth}{-1cm}{}
    \begin{tabular}{>{\centering\arraybackslash}p{6cm}|>{\raggedright\arraybackslash}p{12.5cm}}
      \multirow{5}{6cm}
      % Cột bên trái: ÉPISODE và số tập
      {\\[-40pt]
        \flaregothic\fontsize{20pt}{20pt}\selectfont \centering
        {\contour{errorfour}{\textcolor{black}{ÉPISODE}}}\color{black}\\
        \burbank\fontsize{72pt}{72pt}\selectfont
        \includegraphics[height=3.12359cm]{../../../../Image/Book?/number/num24.png}
        \\[18pt]
        % \flaregothic\fontsize{14pt}{14pt}\selectfont
        % \color{epattr}BIENVENUE, \uline{\color{Banpire} Mel} !
        \flaregothic\fontsize{18pt}{18pt}\selectfont
        {\contour{errorfour}{\textcolor{white}{ENVAHI}}}\\{\contour{errorfour}{\textcolor{white}{FIN DE TOME\dots\ ?}}}
        \color{black}
      }
       &
      % Dòng thứ nhất: tên tập tiếng Nhật
      {\fotskip\fontsize{18pt}{18pt}\selectfont \ruby{青}{あお}\ruby{上}{がみ}\ruby{高}{こう}\ruby{校}{こう}へようこそ、またしても
      } \\[6pt]
      % Dòng thứ hai: tên tập tiếng Anh
       & {\cafeteria\fontsize{18pt}{18pt}\selectfont Welcome to Aogami High, Once Again!}                  \\[4pt]
      % Dòng thứ ba: tên tập tiếng Việt
       & {\vnmsans\fontsize{18pt}{18pt}\selectfont Chào mừng tới Trường Trung học Aogami, một lần nữa!}    \\[8pt]
      % Dòng thứ tư: Nhân vật xuất hiện
       & Track used:
      \begin{itemize}
        \item {\sen{\fontsize{14pt}{14pt}\selectfont \href{https://www.youtube.com/watch?v=Uhu21vKi3As}{Lov-E (\ruby{電}{でん}·\ruby{愛}{あい}) - Kr1TT Remix}}}
        \item {\caviar{\fontsize{14pt}{14pt}\selectfont Dream of Sky (Chamber Remix) - Łukasz Michalski \& 小陳x\_chen}}
      \end{itemize}
      \\[6pt]
      % Dòng cuối cùng: Thời gian diễn ra của tập (thường sẽ là ngày phát hành)
       & {\continuum\fontsize{16pt}{16pt}\selectfont Ngày phát hành: 16/07/2022}                           \\[8pt]
    \end{tabular}
  \end{adjustwidth}
\end{table}
%===========================================NỘI DUNG==========================================
\color{black}

\begin{center}
  \pov{Hikawa Kuon}
  \uline{Nhà Hikawa (tạm) - thị trấn Aogami.\\Tháng Tư, năm 2021\\Năm giờ sáng.}
\end{center}

\dots

``\dots\ Ặc.''

\href{https://www.youtube.com/watch?v=Mqokmn0WOMo}{Tiếng chuông báo thức của điện thoại tôi vang lên.}

Tiếng chuông báo thức quen thuộc vang lên, đều đặn và chói tai. Nhất là vào cái giờ oái oăm thế này.

``Ừm\dots'' tôi khẽ rên, quờ tay tìm chiếc điện thoại vẫn đang rung bần bật đâu đó trên đầu giường.

``\dots\ Ầy da\dots Ay! Đây rồi!''

Ánh sáng từ màn hình lập tức đập vào mắt, khiến tôi phải nheo lại. Năm giờ sáng. Chẳng ai tầm tuổi tôi mà lại dậy sớm thế này - nhưng bằng cách nào đó, tôi đã quen rồi.

Bên cạnh tôi, Kaigou vẫn cuộn tròn trong chăn, mái tóc xoăn bung ra trên gối, nhịp thở đều và nhẹ. Ở phía bên trong cùng, Suzuki đang ngủ say, cánh tay nhỏ vắt ngang qua người chị mình. Nhịp thở của cả hai hoà vào không gian tĩnh lặng, chỉ bị cắt ngang bởi tiếng báo thức mà tôi vừa tắt đi.

Tôi khẽ ngồi dậy, dựa lưng vào thành giường, thở ra một hơi dài. Ngoài kia, màn đêm vẫn chưa tan hết; bầu trời còn lấp lánh vài ngôi sao mờ. Chẳng biết từ khi nào, tôi lại thấy bình yên với cái khung cảnh này đến thế.

\dots

\dots

\dots Khoan đã.

Chẳng phải tôi đã từng thấy cảnh tượng này rồi sao?

Tôi ngồi bất động vài giây, nhìn quanh căn phòng. Căn gác gỗ nhỏ, mùi nhựa thông phảng phất, hai chị em nhà họ vẫn đang say ngủ bên cạnh\dots\ Mọi thứ đều thật, nhưng lại có gì đó trùng khớp một cách kỳ lạ với một ký ức đã bị lãng quên.

Tôi lắc đầu, cố gắng xua đi ý nghĩ ấy.

Từ chỗ tôi ngồi, khung cửa sổ hé mở, để lộ thị trấn Aogami đang dần tỉnh giấc. Không còn là những mái nhà hoang lạnh, không còn là sương mù xám xịt. Thay vào đó là ánh đèn vàng của các ngôi nhà nhỏ, tiếng người gọi nhau chuẩn bị mở hàng, và những tán cây xanh rung rinh trong gió sớm. Thị trấn miền núi ấy - nơi từng ám ảnh tôi trong những giấc mơ lặp đi lặp lại - giờ đang sống lại, thật sự sống.

Tôi tự cười khẽ. Có lẽ vì cái giấc mơ đêm qua.

Một giấc mơ kỳ quái: tôi, Kaigou và Suzuki cùng một cô gái tên Shino đã gieo xuống những hạt giống, vượt qua những thử thách không tưởng, và rồi cả thị trấn Aogami được hồi sinh trong biển sáng trắng. Nó rõ ràng đến mức tôi gần như cảm nhận được cả hơi lạnh trong gió, mùi đất, và giọng nói của ai đó vang lên gọi tên tôi.

``Nghe như mấy phim siêu anh hùng ấy nhỉ\dots'' tôi bật cười nhỏ, rồi vươn vai, chuẩn bị đứng dậy.

Khi ấy, tôi chợt nhận ra có một vật nằm trên bàn: một chiếc đồng hồ đeo tay màu bạc, đang sạc pin. Tôi khựng lại. Mặt đồng hồ hơi sáng lên, phản chiếu ánh bình minh yếu ớt.

Một thoáng rùng mình chạy dọc sống lưng. Chẳng phải đó là chiếc đồng hồ đã cùng tôi chinh chiến trong giấc mơ đó sao?

``Mình đang nghĩ gì chứ\dots\ Tại sao lại tưởng tượng mấy thứ kỳ quái đó vào ngày trọng đại này?'' tôi lắc đầu, cố gắng ngắt mạch suy nghĩ đang dần trôi về phía mơ hồ kia.

Hôm nay là ngày đầu tiên chúng tôi đến trường mới: bạn mới, thầy cô mới, môi trường mới\dots\ mà có lẽ tôi đã quen thuộc tới mức nhàm chán - nếu dựa theo những gì xảy ra trong giấc mơ rồi.

Tôi khẽ thở ra, nhìn ra ngoài cửa sổ thêm lần nữa.

``\textit{Nhưng mà\dots\ có thật là chúng ta đã thoát khỏi vòng lặp đó rồi chứ?}'' tôi thầm nghĩ, nhìn về phía thị trấn đang bắt đầu vận hành ngày mới.

Ánh sáng trắng bao phủ quanh chúng tôi tan dần, để lại cảm giác choáng váng như vừa bước ra khỏi cơn mưa tuyết.

``\textit{Trước tiên thì\dots\ phải gọi hai chị em nhà nó dậy đã,}'' tôi tự nhủ.

Những câu hỏi về vòng lặp và thế giới tạm thời cứ gác sang một bên. Tôi lay vai từng người một, nhẹ giọng gọi.

``Kaigou-chan? Suzuki? Dậy đi nào.''

Cô em út cựa mình, giọng khàn khàn: ``Onii-sama\dots? Đây\dots\ là đâu? Ta có còn\dots\ trong mơ không?''

Cô chị ở bên cạnh lồm cồm bò dậy, ngáp dài, mái tóc rối tung. ``Càm giác đó\dots\ kỳ lạ thật. Như thể mọi thứ trong giấc mơ đó\dots\ có thật vậy.''

``Tôi cũng đang mông lung giống hai người đây,'' tôi thở dài, chống hông nhìn hai ``công chúa ngủ trong rừng'' vừa mới tỉnh dậy sau một giấc ngủ sâu.

Chúng tôi ngồi thành hàng trên sàn, cố ghép lại ký ức trong giấc mơ đó.

Trong giấc mơ, chúng tôi đã cùng nhau gieo hạt xuống nền đất cằn cỗi của thị trấn Aogami, nơi trước đó chỉ còn lại tro tàn và gạch vụn. Tôi thấy từng mầm cây bật lên, từng ngọn cỏ lan ra, rồi đường phố dần trở nên sáng hơn dưới ánh ban mai. Những con đường mà trước đây chúng tôi từng thấy nứt vỡ, nay tràn đầy sức sống.

Rồi chúng tôi gặp hai cô gái.

Một người có đôi mắt xanh trong suốt như pha lê - Misora Shino, tượng trưng của Sora-chan.

Người còn lại, với mái tóc hồng và bộ đồ vu nữ, dịu dàng cười khi nâng tay chào chúng tôi - Shinomiya Sakura. Trông cũng hơi hơi giống Miko-chi.

Cả hai nói chuyện, cùng chúng tôi đi qua những dãy phố hồi sinh, và rồi\dots\ ánh sáng nhấn chìm tất cả.

Kaigou-chan chau mày. ``Nhưng đó là mơ, đúng không? Chúng ta vừa ngủ dậy mà.''

``Em cũng thấy khó tin lắm\dots'' Suzuki ở bên cạnh, gật đầu đồng tình với cô chị. ``Nhưng mà\dots\ em vẫn cảm thấy có gì đó thật sự vừa xảy ra vậy\dots''

Một giọng nói phát ra, đều đặn từ chiếc đồng hồ đeo tay tôi đang cắm sạc trên bàn: ``\eng{Đó là vì thứ các cậu thấy không phải là mơ. Đó là sự thật, một thế giới đã được hồi sinh từ chính ước vọng của các cậu. Cô cậu có để ý rằng thị trấn Aogami này tràn đầy sức sống hơn hẳn cái mà cậu từng thấy trong giấc mơ không?}''

Tôi nhìn ra ngoài cửa sổ. Ánh mặt trời sớm len qua tán cây, rọi xuống mái nhà của những hộ dân bên kia đường. Không còn sương xám, không còn bóng đổ rối loạn hay rạn nứt không gian như trước. Mọi thứ sống động đến mức\dots\ gần như không thể là mơ.

Tôi gật đầu. ``\eng{Đúng là thế thật\dots\ Nhưng còn Shino-san và Sakura-san thì sao? Liệu họ có nhớ gì không?}''

``\eng{Có. Họ tồn tại thật. Và họ nhớ tất cả. Dù họ sẽ không nói ra, ít nhất là cho đến khi các cậu gặp lại.}''

Cả hai chị em nhìn nhau, ánh mắt vừa mông lung vừa ngập ngừng. Chúng tôi im lặng vài giây, trước khi Game tiếp lời, giọng trầm xuống: ``\eng{Giờ thì chuẩn bị quần áo sách vở đi. Hôm nay là ngày đầu tiên các cậu nhập học ở trường Aogami. Đừng để muộn.}''

Cô chị lớn nheo mắt: ``\eng{Sao tự dưng lại gấp thế?}''

Ông ta thở dài. ``\eng{Nhớ những gì Kuon-user, cô và Suzuki-san đã hứa với Shino-san và Sakura-san chứ?}''

Chúng tôi thoáng sững lại.

Tôi khẽ mỉm cười, đáp: ``\eng{Ờ\dots\ đúng rồi. Chúng tôi đã hứa sẽ sống cùng họ, như những học sinh thật sự.}''

Câu trả lời ấy như phép thuật pha loãng bầu không khí nặng nề. Bốn đứa nhìn nhau, mắt ai cũng hơi ánh lên niềm tin mong manh, một thứ không dám nắm chặt, chỉ dám để nó lơ lửng giữa chúng tôi.

Kaigou-chan là người đầu tiên rời khỏi giường. Cô vươn vai, đầu tóc rối tung như tổ quạ, áo ngủ xộc xệch, trông như con mèo vừa trốn khỏi cơn bão. Suzuki nằm co ro thêm năm giây, rồi cũng lê người xuống sàn, mắt nhắm nghiền, miệng lầm bầm: ``Em muốn buổi sáng bị hoãn vô thời hạn\dots''

``Không ai hoãn cả,'' cô ấy cười khẽ, vỗ vai em gái. ``Dậy đi, còn phải làm dáng với đồng phục mới.''

Tôi đã gấp sẵn bộ đồng phục nam, gồm một chiếc áo sơ mi trắng, quần xanh đậm, cà-vạt đồng màu - ngay trên ghế từ tối qua. Còn hai chị em thì lôi bộ nữ từ mắc áo: váy ca-rô xanh thẫm, nơ trước ngực. Suzuki vừa cài nơ đã thở dài: ``Nơ xanh này làm em như bánh bông lan đó, Onee-sama\~{}''

``Nhưng bánh bông lan có ai chê đâu,'' Kaigou đáp, rồi cặp thêm nơ đỏ lên tóc mình, tay khẽ vuốt phần mái lệch. ``Xinh mà.''

Cô em út đỏ ửng tai, quay đi: ``Chị cũng\dots\ không tệ.''

Vì trời còn se lạnh, trường đã phát thêm áo khoác: blazer (áo khoác mỏng) đen và áo cardigan (áo len). Trong ký ức vòng lặp nhà ga, tôi nhớ Kaigou-chan khoác blazer như đội trưởng tháo lui, còn Suzuki co rúm trong cardigan, mũi đỏ hồng, thì thầm: ``Ấm quá, em không muốn cởi ra bao giờ.'' Và bây giờ, vẫn vậy. Không thay đổi gì.

Cả ba chúng tôi lần lượt vào phòng tắm. Kaigou-chan đi trước, mười phút sau mới ra, tóc còn ẩm, mùi dầu gội thoang thoảng. Suzuki ôm khăn, lê bước vào, cố tình để lại cửa hé mở để nghe tiếng chị hát vu vơ. Tôi vào cuối cùng, tranh thủ súc miệng, rồi nhìn gương, thấy mình cũng đang tươi tỉnh hơn mọi ngày.

\ct

Sáu giờ kém năm phút.

Tôi vừa đeo cặp vừa hỏi: ``Sách vở, bút thước, đều ổn chứ?''

``Ổn hết,'' Kaigou đáp, chỉnh balo lên vai. ``Còn cả bánh bao nữa, nếu ai đó đói bụng giữa buổi.''

Suzuki nhăn mũi: ``Em không đói-- à, chỉ hơi thôi.''

``Quên bút thì đừng có mếu máo nhé,'' tôi trêu.

``Em đã lớn rồi!'' cô bé phồng má, nhưng rồi lại thì thầm: ``Nhưng mà\dots\ anh vẫn cho em mượn chứ?''

Kaigou bật cười, xoa đầu em: ``Có cả chị và Kuon-kun làm hậu phương, lo gì.''

Tôi mở cửa, ba chúng tôi bước xuống cầu thang. Tiếng bước chân lúc lỉu hòa vào nhau, như một bản nhạc chưa viết hết, một khúc đầu ngày mới, đủ ấm để chúng tôi không sợ lạc đường.

\ct

Bữa sáng hôm nay diễn ra trong một khung cảnh yên tĩnh đến mức gần như nghe được tiếng đồng hồ tích tắc trong bếp. Mùi gạo vừa nấu chín trộn lẫn với hương trứng rán lan khắp căn phòng nhỏ. Kaigou-chan đang lục tủ, còn Suzuki thì loay hoay chuẩn bị mấy phần rau xào. Cả hai có vẻ thành thạo chuyện bếp núc hơn tôi tưởng.

\dots Mà cũng phải thôi, hai chị em họ đã từng sống nương tựa nhau trong suốt thời gian không còn bố mẹ, nên việc thạo bếp núc là chuyện hiển nhiên mà.

Tôi đứng khoanh tay dựa vào tủ lạnh, đọc mảnh giấy được dán ngay trên cánh cửa:

\txtbx{Hai đứa à, bố mẹ phải đi công cán vài tuần. Tiền sinh hoạt để trên kệ tủ. Ăn uống cho cẩn thận, nhớ nấu ăn đàng hoàng.}

Ngay bên cạnh, một phong bì nhỏ đặt gọn gàng, bên trong là vài tờ tiền còn mới.

Tôi thở dài.

``Cả thế giới này cũng viết kịch bản y chang mấy trò nhập vai ấy nhỉ,'' tôi buông tiếng cười chát, ``bố mẹ bận công tác, để lại tiền, ba anh em tự lo\dots\ nghe quen tới mức phát sợ.''

Kaigou-chan thì chỉ nhún vai, giọng cô nhẹ hẳn đi: ``Ít ra thì yên tĩnh. Không ai hỏi mình học hành tới đâu, hay bắt phải ngủ sớm.''

Suzuki bật cười, tay vẫn đảo chảo trứng: ``Nghe cứ như ba anh chị em đang tận hưởng đời sống tự lập ấy.''

Tôi ngồi xuống bàn, chống cằm. ``Ừ, có khi vậy thật.''

Bữa sáng chỉ có cơm trắng, trứng rán và miso. Nhanh gọn, ấm áp, và đủ năng lượng để bắt đầu ngày mới. Hộp cơm trưa thì khá tạm bợ: ``vài đồ thừa trong tủ lạnh và vài món khác'' như cách tôi mô tả. Nhưng ở thị trấn này, chỉ riêng cảm giác được ngồi ăn trong căn bếp sáng sủa, nghe tiếng gió len qua cửa sổ, cũng đã quá đủ với chúng tôi rồi.

\ct

Trên đường tới trường, tôi khoác cặp sau lưng, tay đeo chiếc đồng hồ vừa được sạc đầy. Hai chị em họ đi song song phía trước. Cả ba im lặng một lúc lâu, chỉ có tiếng bước chân lẫn trong tiếng xe điện xa xa.

Thị trấn Aogami lúc sáu giờ rưỡi sáng mang một vẻ đẹp mà tôi chưa bao giờ chán nhìn, dù là ở trong giấc mơ, hay ở đời thực.

Nhắc đến quang cảnh, nơi này không tấp nập như thành phố, nhưng lại thoáng đãng và sáng sủa. Mặt trời mới lên, ánh nắng len qua tán cây, mùi gỗ ẩm lẫn trong gió mát. Dù đã ở đây một thời gian, mỗi lần đi qua những con phố này, tôi vẫn có cảm giác như lần đầu.

Aogami là một thị trấn nhỏ nằm giữa những dãy núi xanh thẳm, nơi những ngôi nhà truyền thống bằng gỗ vẫn đứng sát cạnh các tòa nhà bê tông. Những con đường lát đá uốn lượn dọc theo sườn núi, dẫn tới những cửa hàng tạp hóa nhỏ xinh và những quán cà phê ấm cúng. Dù không có những trung tâm thương mại hay rạp chiếu phim hoành tráng, nhưng nơi đây vẫn có đầy đủ tiện nghi của cuộc sống hiện đại.

Điều đặc biệt nhất của Aogami có lẽ là cách người dân nơi đây vẫn giữ gìn được nét văn hóa truyền thống giữa dòng chảy của thời đại. Những lễ hội theo mùa vẫn được tổ chức đều đặn, tiếng trống taiko vang vọng từ đền thờ cổ kính trên đỉnh đồi, trong khi những chiếc xe điện chạy êm ru trên những con phố dưới chân núi. Sự hòa quyện giữa cũ và mới tạo nên một bầu không khí thanh bình, khiến bất cứ ai đến đây đều cảm thấy như được trở về nhà.

Tháng Tư ở Aogami vẫn còn đọng lại chút hơi lạnh cuối đông, nhưng những tia nắng xuân đã bắt đầu ấm áp hơn, xuyên qua làn sương sớm còn vương trên đỉnh núi. Gió thổi nhẹ, mang theo hơi se lạnh phảng phất, nhưng không còn căng thẳng như mùa đông nữa. Có thể cảm nhận rõ rệt mùa xuân đang đến gần, qua từng luồng không khí mềm mại hơn, qua những chồi non bắt đầu nhú trên cành cây ven đường.

Càng gần trường, dòng người càng đông.

Học sinh mặc đồng phục bước đi từng nhóm nhỏ, tiếng trò chuyện ríu rít vang khắp ngã tư. Có nhóm con gái cười khúc khích bên quán bánh mì, có nhóm con trai đang đá lon nước ngọt lăn lóc trên vỉa hè.

Kaigou-chan khẽ nói, ánh mắt dõi theo: ``Giống y như trong giấc mơ nhỉ, Kuon-kun?''

Kaigou gật đầu, có vẻ đồng tình: ``Ừ, nhưng lần này\dots\ thật hơn. Cảm giác như mọi người đang thực sự sống, chứ không còn là những cái máy vô hồn nữa.''

Giọng Game vang lên trong đầu tôi, trầm nhưng hơi pha chút thích thú: ``Về cơ bản, có vẻ không khác mấy so với ở trong giấc mơ của các cậu nhỉ?''

Tôi mỉm cười, bước qua cánh cổng sắt cao của Trường Trung học Aogami. ``Khác chứ,'' tôi đáp khẽ, ``lần này, chúng tôi sẽ sống thật.''

Chúng tôi vừa bước qua cổng trường, tiếng chuông đầu tiên trong ngày vang lên, giòn tan như muốn báo hiệu một khởi đầu mới. Dòng học sinh đổ vào sân, áo đồng phục phất phơ trong gió sớm. Mọi thứ bình thường đến mức, gần như tôi phải tự nhéo mình để chắc rằng đây không phải là mơ.

Rồi giữa đám đông ấy, tôi thoáng thấy hai gương mặt thân quen. ``Onii-sama, kia chẳng phải là\dots''

Tôi nhìn theo hướng Suzuki chỉ và suýt nữa thì bật cười.

Hai gương mặt mà tôi tưởng chỉ có thể gặp lại trong mơ, giờ lại đang chí chóe giữa sân trường như chưa từng có chuyện gì xảy ra.

Một cô gái có mái tóc nâu buộc thành hai bím đuôi sam với bờm tai chó đang chống nạnh, trừng mắt với cô bạn tóc tím ngắn bờm tai mèo bên cạnh, đôi má phồng lên như con sóc đang giận vì bị tước mất đồ ăn yêu thích của mình.

``Yuka! Trả hộp cơm lại cho tôi ngay!''

``Không\~{} Bà ăn hết phần của tôi hôm qua rồi còn gì\~{}''

``Đó là vì bà chậm chạp quá chứ bộ! Mồ!''

Akane-san giơ tay định giật lại, trong khi Yuka-san thì tránh sang bên, cười ha hả. Cảnh tượng ``mèo vờn chuột'' ấy sống động đến mức tôi gần như quên rằng, trong một ``thế giới khác'', cô nàng tai chó từng đứng lặng trước cổng trường, lặp đi lặp lại câu ``trả lại bạn cho tôi'' đầy đáng sợ.

Kaigou cười nhẹ, chợt lên tiếng: ``Khác hẳn, nhỉ?''

Tôi gật đầu. ``Ừ. Có vẻ lần này, họ đã tìm lại được nhau thật rồi.''

Nhận ra ánh mắt của chúng tôi, cô nàng tai chó lập tức quay lại, vẫy tay: ``Ơ! Là mấy người hôm mới chuyển đến phải không? Học cùng trường với bọn mình sao?''

Yuka-san nghiêng đầu, nheo mắt: ``Hể, chẳng lẽ cậu là Kuon-senpai mà Akane nói sao? Trông cũng điển trai nha\~{}''

Tôi bật cười: ``Đừng `senpai' kiểu đó, nghe lạ lắm. Gọi Kuon-san thôi được rồi. Cùng khối mà.''

``Thế càng tốt!'' cô bạn tai chó hăng hái. ``Chúng ta học năm nhất, chắc sẽ gặp nhau nhiều đấy!''

``Ừ, hy vọng thế,'' tôi đáp, dù trong lòng vẫn thấy kỳ quặc khi nhắc lại những cái tên đó trong bối cảnh bình thường thế này.

Hai cô nàng lại bắt đầu cãi vã chuyện ai phải mang hộp cơm trưa, rồi vừa cười vừa chạy biến về phía dãy lớp học, hòa vào dòng học sinh đang càng lúc càng đông hơn.

Còn lại chúng tôi đứng ngẩn ra một lúc, trước khi tôi lên tiếng: ``Đi thôi, kẻo trễ. Phải xem lớp lủng thế nào đã.''

\ct

Khu bảng thông báo dán đầy danh sách lớp, tên chen chúc thành từng hàng ngay ngắn. Tôi rà tay tìm từng lớp học, rồi tìm ra tên của mình.

``Hikawa\dots\ Hikawa\dots\ Đâu rồi nhỉ\dots? À, đây rồi.''

Tôi thấy tên mình trong danh sách. Lớp 1-C. Kỳ quái thật, trong giấc mơ, tôi đã thấy mình học ở lớp đó.

Nhưng điều bất ngờ chưa dừng ở đây. Ngay ở phía dưới tên tôi là hai cái tên khác cùng họ với tôi: Hikawa Kaigou và Hikawa Suzuki.

``Vậy ba chúng ta học cùng lớp sao!?'' Suzuki reo lên, mặt sáng bừng.

Tôi cười nhẹ. ``Nghe tiện thật. Không phải lo tìm nhau trong đống người lạ.''

``Kể cả ở trong thực tại thì chúng ta vẫn học chung với nhau nhỉ, Kuon-kun?'' Kaigou-chan húc nhẹ vào vai tôi, cười gian.

``Ừ\dots''

Nhưng khi mắt tôi lướt qua danh sách còn lại, nụ cười bỗng khựng lại. Những cái tên trong này\dots\ không thể nhầm.

Shino-san và Sakura-san, cùng với Akane-san và Yuka-san\dots\ Họ sẽ ở chung lớp với chúng tôi.

``Inari, Hanabi, Touka, Saki, Suzune\dots\ chẳng phải đó là những học sinh từ vòng lặp đầu tiên của chúng ta sao?'' cô chị lớn trố mắt ngạc nhiên khi nhìn thấy những cái tên đó trong lớp 1-C.

``Ừm\dots\ nhưng mà\dots\ vẫn chưa hết đâu. Còn đây này.''

Nanase, Yae, Shiina, Honoka, Mari từ thời quá khứ cũng sẽ ở trong lớp 1-C.

``K-Không đời nào\dots'' cô em út trợn ngược mắt ngạc nhiên. ``Em cũng không ngờ là họ cũng sẽ ở trong lớp của chúng ta đấy.''

Còn thêm nữa: Kaoru, Miku, Saya, Uzuki, Kana, Kanade, Hotaru, Mitsuki, Yuki\dots\ tất cả các thành viên từ vòng lặp nhà ga và từ thế giới ``hoàn hảo'' đó\dots\ họ cũng trong lớp C chúng ta.

``\eng{\dots Xem chừng như chúng ta cũng vô tình làm tái sinh luôn tất cả mọi người nhỉ, Kuon-user,}'' Game-san lên tiếng, giọng nửa quan ngại, nửa buồn cười.

Và sáu cái tên mới toanh, nhưng tôi nhận ra một vài trong số đó từ thế giới gốc của mình: Utachi Azuki, Mizuno Aku, Akaba Koishin, Shitou Shion, Onidoro Ayame, Hitsujikado Wata. Họ sẽ không đền vào hôm nay, nhưng chắc chắn là sẽ chung lớp với chúng tôi.

Tổng cộng, lớp 1-C có ba mươi hai học sinh. Và ngoài tôi ra\dots\ không có lấy một thằng con trai nào khác. \emph{Không một ai, ngoài tôi.}

Tôi nín lặng vài giây. Kaigou-chan nhíu mày, còn Suzuki thì che miệng cười khúc khích. ``Cậu có chắc là cậu không nhập học nhầm vào trường nữ sinh đấy chứ?'' cô chị lớn hỏi, giọng nửa đùa nửa nghi ngại.

Tôi thở dài, mắt nhìn danh sách mà lòng trôi tuột xuống đáy.

\begin{center}
  ``Cuộc đời thật lắm éo le,\\
  Nhân sâm thì ít, mà rễ tre thì nhiều\dots\footnote{Câu lục bát này lấy từ chương trình \textit{Táo Quân} năm 2018.}''
\end{center}

Suzuki phá ra cười, còn Kaigou chỉ lắc đầu, vẻ mặt kiểu ``Kuon-kun lại văn vẻ nữa rồi.'' Nghe những tiếng cười như châm chọc đó lại càng khiến tôi thở dài mạnh hơn, không ngừng tự hỏi vì sao mọi thứ lại sắp xếp có chủ đích như vậy. Tôi biết, việc một người con trai được học cùng với lớp chỉ toàn nữ là điều không hiếm ở thế giới tôi, nhưng thường chỉ gặp ở những lớp theo ban ``học thuộc.'' Tôi chưa bao giờ nghĩ bản thân sẽ rơi vào cảnh hy hữu như thế này.

Bỗng một giọng nói trong trẻo vang lên ngay sau lưng: ``Câu thơ nghe hay đấy, Hikawa-kun. Nhưng hình như cậu sai nhịp rồi.''

Tôi quay lại, và tim như khựng một nhịp.

Shino-san trong bộ đồng phục mới tinh, đang cười nhẹ, mái tóc nâu buông xuống vai. Bên cạnh cô là Sakura-san, vẫn với đôi mắt xanh lá mạ tinh nghịch và nụ cười nửa miệng quen thuộc. Cả hai người họ đang diện đồng phục với áo sơ mi, áo blazer đen khoác ngoài, giống như chúng tôi. Sakura-san thậm chí còn đang mặc áo cardigan màu hồng ở bên trong, khác với màu nâu như bình thường.

``Số thứ tự của cậu trong lớp lại là 13 cơ à?'' Shino nghiêng đầu. ``Mình không biết là nên chúc mừng hay chia buồn nữa.''

``Còn tôi,'' Sakura chống nạnh, ``thì thấy đúng kiểu mơ ước của mấy cậu học sinh đấy chứ? Một mình giữa cả rừng hoa thế kia mà.''

Tôi cười gượng. ``Phải, chỉ là\dots\ tôi không nghĩ giấc mơ lại trở thành hiện thực tới mức như thế này.''

Kaigou khoanh tay, chen ngang, giọng đầy tự tin: ``Nếu bất kỳ tên nào dám bén mảng tới cậu, tôi sẽ cho hắn về lại kiếp mơ luôn.''

Suzuki đỏ mặt, lí nhí: ``E-em sẽ\dots\ cổ vũ anh.''

Tôi bật cười, xoa đầu con bé. ``Anh đâu cần đội bảo vệ đâu mà mạnh ghê vậy.''

Shino mỉm cười dịu dàng, còn Sakura chỉ nhún vai: ``Vậy thì tốt. Đi thôi, lớp 1-C ở dãy phía Tây đó.''

Tôi bước sát lại, vai chạm vai hai cô gái; hơi ấm từ cánh tay họ truyền qua lớp áo đồng phục mỏng, như thể muốn xác nhận rằng tất cả đã thật sự trở lại. Trong khoảnh khắc ấy, tiếng chuông báo hiệu giờ học vang lên đúng một hồi, gió sân trường cuốn theo mùi hoa anh đào thoang thoảng, và tôi chợt nhận ra mình đang thở đều, chậm, trọn vẹn, điều mà suốt chuỗi ngày mơ hồ trước đó tôi chưa từng làm được.

Ánh mắt Shino long lanh như hồ nước thu, khóe môi cong lên một cung tròn mềm mại không chỉ là niềm vui, mà còn là cả một bầu trời nhẹ nhõm sau bao ngày chờ đợi. Gò má cô ửng hồng, những sợi tóc mai bồng bềnh rung nhẹ trong gió sớm, phản chiếu nắng thành từng tia vàng li ti, như thể chính ánh nắng cũng reo vì được thấy cô cười. Đôi tay cô đan trước ngực, ngón tay luồn vào nhau chặt hơn một chút, như vừa tìm lại được thứ quý giá nhất.

Sakura bên cạnh tuy chỉ nhún vai, nhưng đuôi mắt cô hếch lên, lộ ra tia sáng ẩn như pha lê đang được rót đầy; cánh mũi nhẹ phập phồng, bật ra tiếng thở dài ngắn cũn cỡn như tiếng ``hừ'' khe khẽ, chất chứa niềm hạnh phúc không thể che giấu. Cả hai không nói gì thêm, nhưng không gian quanh họ như rung lên một nhịp, gọi tên sự sống mới.

Và thế là, giữa dòng học sinh rộn ràng của buổi sáng đầu tiên, năm người chúng tôi cùng bước vào lớp học mới, nơi mà giấc mơ và hiện thực, có lẽ, cuối cùng đã hòa làm một.

\ct

Chúng tôi rảo bước dọc theo hành lang, vừa tìm lớp vừa ngó nghiêng xung quanh. Khu của khối nhất nằm ở dãy phía Tây, như Sakura-san vừa nói.

Nhưng mà\dots\ không hiểu sao tôi lại thấy có chút là lạ. Tôi đang nói đến cái hành lang này.

Sàn gạch màu ngà, trần kẻ ô, mấy cái loa phát thanh được gắn đều tăm tắp, bên phải là dãy cửa sổ lớn hắt ánh sáng trắng xóa xuống nền gạch, bên trái là những lớp học có bảng thông báo màu xanh, dán đầy giấy và ghi chú chi chít.

Trước mỗi cửa lớp, trên bậu cửa sổ hoặc chiếc bàn nhỏ góc hành lang, đặt những chậu hoa bỉ ngạn \textbf{trắng} tinh khôi, cánh mỏng manh như sương sớm, đung đưa nhẹ trong làn gió lọt qua khung kính.

\img{e3-8a}

Tôi khựng lại một giây. ``Bỉ ngạn trắng ư?''

Trong mơ, tôi chỉ từng thấy bỉ ngạn đỏ rực - là màu của chia ly, của tử biệt. Vậy mà tại sao lại\dots

Game-san lập tức lên tiếng, như thể đọc được suy nghĩ của tôi: ``\eng{Đỏ là dấu hiệu kết thúc, còn trắng\dots\ là khởi đầu.}''

Tôi nhướng mày. ``\eng{Thật ư?}''

Giọng ông ta lúc này trầm xuống, như thể đang đọc một đoạn mã nguồn bị lãng quên. ``\eng{Thuần khiết, phục sinh, tái sinh sau cái chết - loài hoa ấy thích hợp cho những linh hồn vừa được khởi động trở lại.}''

Tôi nhìn theo bóng Shino và Sakura đang bước phía trước, váy đồng phục phất phơ. Họ từng là ``người trong mơ'', nay lại hiện hữu nơi hành lang đầy nắng này - đúng như lời Game nói: \textit{tái sinh}. Lúc này, tôi mới hiểu vì sao người ta đã đặt trước cửa lớp những chậu bỉ ngạn trắng, như một lời nhắc khẽ: \textit{Chào mừng trở lại}.

Tất cả mọi thứ không khác đi một ly so với trong trí nhớ của tôi, hay đúng hơn là trong những giấc mơ lặp đi lặp lại ấy. Chỉ có điều, trông nơi này tươm tất và sạch sẽ hơn nhiều. Không còn vết nứt trên tường, không còn cảm giác ẩm mốc và lạnh lẽo. Mọi thứ đều sáng rõ, sống động, đầy ắp hơi người.

Cảm giác thật đến mức tôi chẳng dám nghĩ là mình mơ nữa.

Bình tĩnh lại nào, Hikawa Kuon. Đây là thực tại: một thế giới thật, một cuộc sống thật. Mọi thứ trùng khớp kia\dots\ chỉ là ngẫu nhiên thôi.

Tôi tự nhủ như thế, dù trong lòng vẫn còn vương lại chút nặng trĩu.

Hy vọng rằng ngày đầu tiên ở đây sẽ diễn ra suôn sẻ\dots\ ngoại trừ việc bốn phương tám hướng đều là con gái cả. Chậc.

\ct

Khi chúng tôi bước vào lớp 1-C, cả năm người đều hơi khựng lại trước bầu không khí trong lành và náo nhiệt của buổi sáng đầu tiên.

Phòng học sáng trưng, bàn ghế ngay ngắn, mùi gỗ mới pha lẫn hương hoa từ sân trường thổi vào qua khung cửa sổ mở toang.

Giáo viên chủ nhiệm tạm thời - một cô giáo trẻ có giọng nói dễ chịu, giới thiệu rằng cô chỉ phụ trách buổi đầu, vì ``giáo viên chính'' sẽ tới vào ngày mai. Cô mỉm cười dịu dàng và bắt đầu phần giới thiệu từng học sinh.

``Nào mọi người, chúng ta hãy cùng nhau giới thiệu bản thân mình trước cả lớp nào!''

Shino giới thiệu đầu tiên. Nụ cười của cô ấy sáng rực, khác hẳn vẻ lạnh nhạt, trầm lặng trong giấc mơ trước đây. Có thể thấy rõ trong ánh mắt ấy là niềm vui thuần khiết, niềm vui của một người thật sự đang sống.

Sau đó đến lượt Kaigou-chan. Cô đứng dậy, vuốt phẳng tà váy, giọng vang rõ ràng: ``Em là Hikawa Kaigou. Rất vui được làm quen với cả lớp.''

Cô vừa dứt lời, Touka ở bàn bên đã nghiêng đầu, mắt sáng lên: ``Khoan đã, hai bạn họ Hikawa này\dots\ mắt và góc hàm giống nhau quá. Không lẽ\dots\ các cậu anh em sinh đôi?''

Tiếng xì xào bùng lên. ``Ồ, đúng nha, nhìn kỹ thì giống thật!'' ``Nhưng một trai một gái á?'' ``Hiếm đó chứ!''

Kaigou-chan chỉ cười nhẹ, lắc đầu: ``Không phải đâu. Kuon-kun và em chỉ là cùng họ thôi.'' Cô quay sang tôi, mắt nheo lại đầy ẩn ý. ``Nhà em có truyền thống nhận con nuôi từ nhỏ; Kuon-kun được bố mẹ em nhận về khi em lên năm, nên chúng em lớn lên bên nhau. Vì thế giống nhau một chút cũng\dots\ bình thường mà.''

Cả lớp ``ồ'' lên như vừa nghe một bí mật vừa được phát hiện. Tôi thì giật mình: chưa bao giờ nghe cô bịa kịch nhanh đến vậy! Nhưng nhìn ánh mắt cô lúc này - vừa thách thức vừa cầu cứu - tôi chỉ còn biết gật đầu bổ sung: ``Đúng rồi, em cảm thấy mình may mắn lắm khi được làm\dots\ anh trai cả của Kaigou-chan và Suzuki-chan.''

Touka-san gãi má, có vẻ hơi tiếc: ``Ra vậy\dots\ nhưng mà nhìn vẫn giống thật.''

Kaigou-chan cúi đầu lịch sự: ``Nếu sau này cậu cần mẫu ADN, cứ gọi bọn tôi là được rồi.'' Cả lớp bật cười, không khí giảm căng ngay lập tức. Cô ngồi xuống, khẽ thở phào, tay vẫy vẫy quạt mát cho mình. Tôi liếc sang, cô nhoẻn miệng thầm: ``\vie{\hl{Kịch bản tạm ổn chứ?}}'' Tôi chỉ kịp cười trừ: ``\vie{\hl{Cứ cho là vậy đi.}}''

Sau đó là phần giới thiệu của những học sinh khác. Tôi cố ghi nhớ từng cái tên: Nanase, Kaoru, Saya, Hotaru, Inari, Yae, và nhiều người nữa.

Nanase thì ngay thẳng đến mức lời nói nào cũng như lời thề. Kaoru có phong thái dứt khoát, giọng nói mạnh mẽ. Saya thì nhẹ nhàng, từ tốn, giọng nói ngân ra như có chút âm nhạc trong đó. Hotaru thì hoạt bát, thân thiện, nhưng có gì đó\dots\ hơi ``lắm mồm'' kiểu người khiến lớp học lúc nào cũng rộn ràng.

Ngồi nghe họ lần lượt giới thiệu, tôi chợt nhận ra, dù không còn trong mơ, tôi vẫn cảm thấy quen thuộc với họ. Một cảm giác vừa ấm áp, vừa lạ lẫm. Giống như đã từng sống cùng nhau, từng trải qua điều gì đó sâu đậm\dots\ nhưng chẳng ai trong số họ nhớ nổi.

Tôi khẽ hít vào, nhìn ra cửa sổ, ánh nắng tháng Tư vẫn chiếu nghiêng qua rèm, phản chiếu lên gương mặt của Kaigou và Suzuki.

``\hl{Thế giới này\dots{} là thực rồi đấy, nhỉ?}'' tôi lẩm bẩm, chỉ đủ để mình nghe thấy.

\pov{Hikawa Kaigou}

Ba tiết học buổi sáng trôi qua nhanh hơn tôi tưởng. Tiết đầu tiên là Văn, môn sở trường của Suzu. Em ấy ngồi thẳng lưng, chăm chú nhìn lên bảng, thỉnh thoảng còn giơ tay trả lời trước cả khi cô giáo kịp dứt câu hỏi. Mỗi lần Suzu nói, giọng vang lên nhẹ nhưng rành rọt, tự tin đến mức khiến cả lớp phải ngoái lại nhìn. Tôi không ngạc nhiên, vì từ hồi còn ở ``thế giới kia'', Suzu đã luôn thích đọc sách, và những bài viết của em ấy thường rất giàu cảm xúc.

Còn tôi thì\dots\ may mà chỉ hơi buồn ngủ chút xíu. Mấy tiết kiểu này thường khiến đầu óc tôi như đang trôi bồng bềnh trong gió xuân. Còn đỡ hơn Kuon-kun; cậu ta thì gục hẳn xuống bàn luôn, chẳng để ý gì tới trời đất nữa.

Sang tiết thứ hai, Vật lý. À, phần này thì ngược lại. Tôi gần như nắm bắt được từng công thức, từng ví dụ. Có vài đoạn tôi còn tự tính nhanh hơn cả thầy, nhưng tất nhiên, tôi không muốn thể hiện quá lố. Kuon-kun thỉnh thoảng liếc sang, khẽ gật đầu như ngầm bảo ``không tệ'', còn Suzu thì hí hoáy ghi chép không ngừng.

Rồi tới tiết thứ ba, Toán đại số, lãnh địa của cậu con trai duy nhất của lớp.

Cậu ta trông hệt như một con cá được quăng lại xuống nước. Tay cầm phấn viết lên bảng loằng ngoằng, giải luôn cả ví dụ khi thầy chưa kịp làm xong. Nhiều bạn gái trong lớp nhìn theo, vừa khâm phục vừa\dots\ có phần hoảng sợ. Tôi biết cậu ta đã có hai năm học Đại học ngành Khoa học máy tính, và ngành đó thì dày đặc công thức và tính toán, nên những gì cậu ta thể hiện đôi lúc lại\dots\ quá trớn.

Còn tôi? Tôi thì chỉ biết thở dài. ``Cái tên này đúng là\dots\ chỉ có thể sống bằng công thức.''

Đến giữa tiết, tôi thấy rõ sự khác biệt giữa ``ba kiểu người'': Kuon-kun và vài người như Nanase-san, Sakura-san thì hoàn toàn tập trung; Suzu và Shino-san thì cố gắng ghi chú tỉ mỉ từng dòng; còn phần còn lại của lớp thì đã bắt đầu rệu rã. Càng về cuối, những ánh mắt ngơ ngác, những tiếng thở dài, và cả mấy cái ngáp kéo dài bắt đầu xuất hiện dọc theo hàng ghế.

Rồi chuông reo, đồng nghĩa với sự giải thoát. Dù chỉ là trong mười lăm phút ngắn ngủi. Ngay khi tiếng ``leng keng'' cuối cùng vừa dứt, thầy dạy Toán cũng gấp sổ bước ra khỏi lớp. Một loạt tiếng vươn vai, tiếng ghế dịch và tiếng trò chuyện vang lên.

Yae-san là người đầu tiên lên tiếng, vừa đi về phía bàn Honoka-san vừa thở phào. ``Phù, cuối cùng cũng kết thúc rồi nhỉ.''

Cô bạn tai Cáo sa mạc thì nằm vật ra bàn, tay vẫy vẫy yếu ớt. ``Mệt quớ\~{} Tiết Toán vừa rồi khiến tớ thấy nhức đầu luôn á\~{} Toàn thấy hàm số với phương trình gì vậy trời\~{}''

\img{e3-8b}

Tôi bật cười gượng, chống cằm. ``Thật ra nó cũng không đến mức khó hiểu lắm đâu\dots\ nếu cậu không nhầm giữa biến số và hằng số thôi. Mà, tôi cũng chẳng khá hơn, Toán đại đâu phải sở trường của tôi; tôi hợp với hình học hơn.''

Shiina-san ngồi ở bàn bên cạnh, gãi đầu, cười chát. ``Có đứa kỳ quái nào lại thích học mấy thứ này nhỉ? Tôi cảm thấy nó chẳng giúp ích gì mấy\dots''

Yae-san gật đầu, vẻ đồng cảm. ``Tớ cũng thấy thế\dots\ Nhưng mà, chúng ta vẫn phải học thôi. Với lại\dots\ công nhận, bài này hơi quá sức thật.''

Honoka-san bật dậy, cau mày. ``Tớ thì không trụ nổi mà cuối cùng ngủ gật tới cuối luôn đây này!''

``Bài học càng khó thì càng làm cho chúng ta buồn ngủ hơn mà,'' cô bạn sư tử vừa nói vừa ngáp dài.

Tôi liếc sang bên kia. Kuon-kun vẫn ngồi nghiêng người trên bàn, tay xoay bút, mắt dán vào tờ bài tập. Cậu ta đã giải gần hết các bài trong sách, lại còn ghi thêm cả mấy biến thể khác mà chẳng ai yêu cầu.

Cô bạn Thể lực trố mắt, thì thào với cô bạn tai cáo Sa mạc: ``Cậu ấy\dots\ giải nhanh hơn cả thầy giáo nữa!'' Honoka-san gật mạnh, hai má phồng lên: ``Như kiểu máy in ra đáp án ấy!''

Cả hai không kìm được mà cùng đồng thanh: ``Khủng khiếp quá đi!''

Tôi mỉm cười gian, quay lại với nhóm con gái. ``Nếu các cậu thấy không hiểu, sao không đi hỏi Kuon-kun? Cậu ta ghiền mấy cái môn này lắm. Với cả, học thầy không tày học bạn mà, đúng chứ?''

Ba khuôn mặt đồng loạt hướng theo ánh nhìn của tôi, rồi cùng gật đầu như thể vừa được khai sáng.

Chỉ vài giây sau, cả ba người - Yae-san, Honoka-san và Shiina-san - đã lũ lượt kéo ghế, tiến về phía bàn Kuon-kun. Tôi đi theo phía sau, vừa cười vừa khoanh tay.

``Kuon-kun\~{},'' cô bạn Thể lực lên tiếng trước, giọng ngọt hơn bình thường, ``cậu chỉ bọn tớ cái chỗ này với\dots\ công thức đạo hàm bậc hai này tớ không hiểu.''

Cô bạn cáo Sa mạc chen ngang, chìa vở ra luôn: ``Cả cái này nữa, tớ ghi được nửa dòng là đầu quay như chong chóng rồi!''

Còn cô bạn Sư tử thì gãi đầu, hạ giọng: ``Tôi\dots\ tôi quên mất quy tắc nhân rồi.''

Kuon-kun ngẩng lên, nhìn ba cô gái đang dần sán tới chỗ cậu, ánh mắt như thể không biết nên cười hay nên khóc. ``Tự dưng lại đi vồ mồi tôi là có ý gì đây?''

``Ý là đang nhờ vả cậu đấy,'' tôi xen vào, nhướng mày trêu. ``Cậu không định sẽ từ chối lời thỉnh cậu của họ đấy chứ?''

Cậu thở dài, đầu hàng: ``Thôi được, đưa đây. Tôi giải cho, nhưng nếu lần sau còn quên, thì tự giác cúng tiền tôi đi nộp học phí đấy nha.''

Chúng tôi cùng cười. Cảnh tượng đó, trong khoảnh khắc, thật bình dị, như thể chưa từng có vòng lặp, chưa từng có những ký ức đau thương nào tồn tại cả.

\pov{Hikawa Suzuki}

Tôi đứng cạnh Shino ở góc lớp, tay ôm tập vở, nhìn về phía trước, nơi Onii-sama đang bị ``bao vây'' bởi một nhóm con gái háo hức hỏi bài, trong khi Onee-sama thì tiếp tục đốc thúc họ. Trông cảnh tượng ấy, tôi không biết nên thương cho anh ấy hay nên buồn cười nữa.

Onii-sama vẫn điềm nhiên như thường lệ, cố gắng giải thích cho từng người, trong khi Honoka-san và Yae-san thì liên tục chen ngang bằng đủ kiểu câu hỏi. Còn Onee-sama thì chỉ đứng khoanh tay phía sau, cười khoái chí theo kiểu ``ta đã nói mà'', rõ ràng đang thích thú khi thấy anh bị ``vùi dập'' trong cơn mưa công thức.

Tôi và Shino-san nhìn nhau, cùng nở một nụ cười mỉm. ``Trông như Onii-sama đang bị vây trong chiến trường vậy,'' tôi khẽ nói.

Cô ấy cười khẽ: ``Ừ\dots\ nhưng ít ra là cậu ấy trông thấy vui. Mình không nghĩ là sẽ được trông thấy cảnh đó.''

Tôi gật đầu, mắt dõi theo ánh sáng đang hắt qua khung cửa sổ. Từ khi bước vào thế giới này, mọi thứ đều có vẻ bình thường đến lạ thường: tiếng gió, tiếng nói cười, thậm chí cả tiếng bút chạm giấy cũng thật hơn bao giờ hết.

``Shino-san này,'' tôi quay sang hỏi, ``cảm giác thế nào\dots\ khi được sống ở thế giới thật?''

Cô ấy ngước nhìn trần nhà, mỉm cười nhẹ.

``Giống như\dots\ ước mơ trở thành thật vậy. Mình đã từng nghĩ rằng mình sẽ mãi bị mắc kẹt trong vòng lặp đó. Nhưng giờ\dots'' cô nắm nhẹ tay áo, ``\dots mình có thể thở, có thể cảm nhận gió, và có thể nhìn thấy họ mỉm cười. Chỉ thế thôi, nhưng với mình, thế là đủ rồi.''

Tôi nhìn cô ấy, định nói gì đó, thì một giọng nói khác bất ngờ vang lên từ cửa lớp: ``\dots Này.''

Giọng nữ có chút kiêu kỳ, rõ ràng và dứt khoát. Shino-san giật mình, suýt đánh rơi bút. Cả hai chúng tôi quay đầu lại. Furukawa Nanase đang đứng đó, một tay đặt ngang chống lên tay còn lại, khuôn mặt lạnh tanh như thường lệ.

\img{e3-8c}

``Eek! C-Chuyện gì vậy\dots?'' cậu ấy lắp bắp.

``Đừng có thù lù xuất hiện như vậy chứ, Nanase-san\dots'' tôi cười gượng.

Cô nàng tóc xanh nhướng mày, giọng điềm đạm: ``Tôi không có cắn các cậu đâu mà sợ.''

Câu nói đó làm không khí dịu đi đôi chút, rồi cô hắng giọng. ``Dù gì thì\dots\ tôi muốn hỏi hai cậu một chuyện.''

``Chuyện\dots\ gì thế?'' Shino-san hỏi lại, vẫn có chút cảnh giác.

Nanase-san khoanh tay, nhìn thẳng: ``Liệu hai cậu có muốn đi xem phim với các bạn khác vào cuối tuần này không?''

``C-Chúng mình á!?'' cả tôi và Shino-san đồng thanh, tự chỉ tay vào mình.

``Chứ còn ai nữa?'' cô ấy đáp gọn, rồi bỗng nở một nụ cười hiền, trái hẳn với gương mặt nghiêm túc trước đó. ``Hai cậu vừa mới vào học mà, đúng không? Tôi chỉ muốn làm quen thêm thôi.''

Tôi thoáng ngạc nhiên. Trong trí nhớ mơ hồ về ``vòng lặp kia'', Nanase-san từng là người lạnh lùng, khép kín và khó gần. Nhưng giờ, ánh mắt cô lại sáng và chân thành đến lạ.

Shino-san mừng rỡ gật đầu: ``Được chứ! Mình cũng rất muốn đi!''

Tôi cũng cười theo: ``Nghe thú vị đấy.''

``Tất nhiên các cậu có thể rủ thêm người khác, tôi không ngại đâu,'' cô bạn tóc xanh tiếp tục. Rồi ánh mắt cô khựng lại. ``Hikawa Kuon-san thì tôi không có ý kiến, nhưng Hikawa Kaigou-san thì\dots''

Cô chợt tái mặt, trông như vừa nhớ lại một điều gì đó kinh khủng. Tôi nghiêng đầu: ``Có chuyện gì sao, Nanase-san?''

Cô lắc đầu, tránh ánh nhìn của tôi. ``K-Không có gì đâu. Chỉ là\dots\ tự dưng tôi thấy chị của cậu\dots\ hơi đáng sợ.''

Tôi bật cười nhỏ, đáp khẽ: ``Onee-sama của mình có thể hơi nghiêm khắc thật, nhưng chị ấy không phải người xấu đâu. Nếu chị ấy có khiến cậu sợ, thì mình nghĩ\dots\ chị đang bảo vệ ai đó thôi.''

Nanase-san thoáng sững lại, rồi mím môi, khẽ nói: ``Có lẽ vậy.'' Cô hắng giọng một lần nữa, cố lấy lại vẻ bình thường.

``Thôi được rồi, vậy thống nhất nhé. Bọn mình đi xem phim cuối tuần.''

Tôi gật đầu, còn Shino thì gần như reo lên: ``Mình mong chờ quá luôn!''

``À mà, Shino-san này,'' con gái của gia đình Furukawa chợt gọi lại, ``cậu có đang bận việc gì không đấy? Tôi thấy cậu định đi đâu thì phải.''

``À, đúng rồi!'' Shino-san khẽ đập tay vào trán. ``Mình định đi làm quen với các bạn khác mà quên mất.''

``Vậy cậu đi đi,'' tôi nói, ``mình sẽ ra sau. Cậu cứ tự nhiên trước nhé.''

``Vậy thì\dots\ mình xin phép. Chốc nữa gặp lại nhé?'' Shino-san mỉm cười, rồi chạy ra ngoài lớp, mái tóc cô khẽ bay trong luồng gió từ hành lang.

Sau khi cậu ấy đi, cô bạn tóc xanh cũng quay lại bàn mình. Nhưng ngay khi cô ngồi xuống, Shiina-san và Honoka-san - ba người vừa hỏi bài từ Onii-sama - tiếp cận tới gần cô ấy.

Cô bạn Sư tử nhướng mày, nheo mắt: ``Chuyện lạ nha. Tự dưng cậu lại chủ động rủ ai đó đi chơi á?''

Nanase-san không đáp. Cô chỉ nhìn ra phía cửa trượt, nơi Shino-san vừa rời đi, ánh mắt như chìm vào suy nghĩ.

``Sao vậy?'' cô bạn Cáo sa mạc hỏi.

Cô bạn tóc xanh khẽ thở ra: ``Kỳ lạ thật đấy\dots\ tôi cảm giác như chuyện này\dots\ đã từng xảy ra ở đâu rồi.''

``Cái đó gọi là déjà vu đó!'' Honoka-san leo lẻo đáp, trong khi Shiina-san chỉ tặc lưỡi: ``Không ai hỏi mà Bộ trưởng trả lời đâu.''

Tôi khẽ bật cười. Không khí trở lại bình thường, nhưng trong lòng tôi lại dấy lên một cảm giác lạ. Quả thật, đúng là Nanase-san đã từng hỏi Shino-san và chúng tôi như vậy ở thời quá khứ, nhưng rồi\dots

\dots một câu hỏi khác bỗng chốc nảy lên trong đầu. ``\textit{Liệu mọi người\dots\ có còn nhớ điều gì từ quá khứ không nhỉ?}''

``Tôi nghĩ vậy, rồi tự lắc đầu. \textit{Thôi thì, cứ để mọi thứ yên ổn như bây giờ đi\dots}''

\ct

Tôi, cùng với Onee-sama và Shino-san, thì ra ngoài hít thở không khí trong lành, và tiện thể, tiếp tục làm quen thêm các bạn học khác. Trong khi đó, Onii-sama vẫn bị vây kín ở trong lớp, hết nhóm này tới nhóm khác thay nhau hỏi bài, như thể cả lớp đều quyết tâm tận dụng hết ``máy giải đề sống'' duy nhất còn sót lại.

Hành lang sáng hẳn lên bởi ánh mặt trời chiếu xiên qua khung cửa sổ. Gió buổi trưa lùa nhẹ, mang theo mùi nắng, mùi phấn trắng, và cả âm thanh xì xào của những lớp bên cạnh.

``Kuon-san đúng là nổi tiếng thật đấy nhỉ,'' Shino-san bật cười, ánh mắt cô vẫn hướng về phía lớp học phía sau lưng.

Tôi cũng không nhịn được mà cười chát, đáp: ``Có lẽ bởi vì trong lớp mình có mỗi anh ấy là con trai thôi\dots\ Mà nói thật thì, nhìn cảnh đó, em cũng thấy tội cho anh ấy ghê.''

Onee-sama khoanh tay, cười xòa, giọng nhẹ tênh: ``Dĩ nhiên chị biết. Nhưng mà, chị cũng đâu thể phụ giúp mấy bài toán kiểu đó được, chị có mạnh môn Đại số đâu.''

Tôi liền quay sang trêu: ``Thế Onee-sama không tính giúp anh ấy một tay à?''

Chị tôi chỉ bật cười, xua tay: ``Cậu ta không phải môn nào cũng giỏi mà. Chị để cậu ta tập trung vào điểm mạnh, rồi bọn chị lo phần khác, chia nhau cho dễ.''

Shino-san gật đầu, tán đồng: ``Đúng rồi, ai cũng có thế mạnh riêng của mình mà. Nếu cùng hỗ trợ lẫn nhau, chắc chắn kết quả sẽ tốt hơn.''

Cả ba chúng tôi bật cười nhẹ, vừa nói vừa đi dọc hành lang. Lúc đó, từ phía xa đầu bên kia, ba cái bóng khác xuất hiện - Hanabi-san, Inari-san và Suzune-san, nếu tôi nhớ không nhầm.

\img{e3-8d}

Cô bạn tóc nâu vừa đi vừa ngáp ngắn ngáp dài, tóc hơi rối, trông rõ vẻ mệt mỏi: ``Trời ạ\~{}\dots\ Mới có buổi học đầu tiên thôi mà sao tôi thấy dài như cả thế kỷ vậy á\~{}''

Cô bạn bờm cáo trắng thì điềm đạm hơn, giọng nhỏ nhẹ mà mang chút ngạc nhiên: ``Bà nói quá rồi, Hanabi-chan. Tôi thấy cũng đâu có đến nỗi đó.''

Còn cô bạn tóc vàng đang đi kế bên thì chợt giơ tay chỉ về phía bọn tôi, mắt sáng lên như đèn pha: ``Ơ, nhìn kìa nhìn kìa! Là Shino-chan với hai chị em nhà Hikawa đó!''

Cô hô to đến mức cả dãy hành lang giật mình. ``\textbf{Này\~{}! Đằng ấy có chuyện gì thế\~{}!?}''

Inari lập tức hoảng hốt, mặt đỏ lên: ``N-Này! Đừng có hét ở hành lang vậy chứ, Suzune-san!''

Nghe thấy tiếng hét, Shino-san khựng lại, nhưng rồi cũng bật cười, vẫy tay: ``A! Mọi người!''

Tôi thì chậm rãi hơn, chào đúng phép: ``Suzune-san, Inari-san, Hanabi-san. Ờm\dots\ đúng không nhỉ?''

Hanabi-san cười tươi: ``Chuẩn luôn! Là bọn mình đó.''

Suzune-san thì gãi má, có vẻ xấu hổ vì màn hô hoán có phần hơi quá vừa rồi.

Còn Inari-san thì lúng túng hẳn, giọng run run: ``À, a-a-nô\dots\ Nếu các cậu bị lạc, thì b-bọn mình có thể chỉ đường giúp\dots''

Onee-sama bật cười ngặt nghẽo, khoanh tay trước ngực: ``Không, không cần đâu. Bọn tôi biết lớp mình ở đâu rồi mà.''

Shino cuống quýt gật đầu lia lịa, đỏ mặt: ``Đ-Đúng! Chúng mình ra ngoài chỉ để hóng gió thôi ấy mà! Không sao đâu!''

Hanabi-san quay sang cô bạn thân (tôi đoán vậy), lắc đầu ngán ngẩm: ``Đấy, thấy chưa, Inari? Bà làm người ta bối rối rồi đó. Bình tĩnh chút đi.''

Suzune vẫn còn đang nhìn quanh, rồi hỏi với vẻ hồn nhiên: ``À-rế, Kuon-kun đâu rồi?''

Onee-sama liếc về phía lớp học, đáp lại với giọng nửa đùa nửa thật: ``Đang làm nhiệm vụ cao cả trong lớp kìa.''

Ba cô bạn đồng loạt nghiêng đầu ngơ ngác.

Tôi đành phải lên tiếng giải thích: ``Ý của Onee-sama là anh ấy đang giúp các bạn khác làm bài tập đấy. Mặc dù\dots'' tôi thở ra một tiếng khẽ, ``\dots mình nghĩ anh ấy bị dồn vào thế không thể từ chối thì đúng hơn.''

Inari-san nghe vậy, khẽ nghiêng đầu, vẻ hơi thương cảm: ``Ừm\dots\ mình hiểu cảm giác đó. Ai nhờ mà từ chối thì kỳ quá, mà giúp hết thì mệt lắm nhỉ\dots''

Suzune thì bật cười, hai mắt sáng rỡ: ``Thế mới bảo Kuon-kun là `hiệp sĩ bàn học' đó! Hôm nay có khi cậu ta nổi tiếng hơn cả giáo viên mất thôi!''

Hanabi khoanh tay, giọng dứt khoát: ``Ờ, giỏi thì giỏi thật, nhưng mà nếu là tôi thì tôi không chịu nổi đâu. Cậu ấy cũng không phải là người duy nhất giỏi Toán đại số trong lớp mình đâu mà.''

Tôi chỉ cười nhẹ. ``Đúng là Onii-sama mà,'' tôi nghĩ thầm, dù có ném anh ấy vào giữa chiến trường nào đi nữa, Onii-sama vẫn sẽ làm mọi thứ trôi chảy như thể đó là việc hiển nhiên, ắt cũng vì nếu tôi nhớ không nhầm, anh ấy cũng là một tổng quản của một công ty giải trí mà.

\ct

\pov{Hikawa Kaigou}

Sau khi chào tạm biệt Hanabi, Inari và Suzune - ba người họ muốn vào lớp sớm để chuẩn bị sách vở cho tiết tiếp theo - tôi, Suzu và Shino vẫn nán lại ngoài hành lang.

Một cơn gió xuân nhẹ khẽ lùa qua, làm tà váy đồng phục của chúng tôi khẽ lay động. Những cánh hoa anh đào từ sân trường bay lả tả qua cửa sổ mở, xoay vòng giữa không trung rồi đáp xuống nền gạch sáng loáng. Mùi hương thanh dịu thoảng trong không khí, dịu dàng, ấm áp, đúng như tháng Tư ở Nhật Bản - theo những gì mà tôi biết.

\img{e3-8e}

Tôi ngẩng nhìn qua khung cửa. Từng cánh hoa hồng phấn rơi xuống như mưa, phản chiếu ánh nắng lấp lánh. Đẹp đến mức khiến người ta muốn nín thở.

``Cậu thấy hoa anh đào đẹp chứ?'' một giọng nói trong trẻo vang lên, pha chút tự hào.

Tôi quay đầu lại. Đó là Ibari Saki-san, cô gái tóc vàng cột hai bên, đang chắp hai tay trước ngực. Ánh mắt cô sáng rực, như thể đang chờ được ai đó đồng tình với mình.

Shino-san hơi lúng túng, đáp lại trong vô thức: ``À, ờm\dots\ Ừ, đẹp lắm.''

Thế là đủ để cô ấy bắt đầu ``bài giảng'' của đời mình. ``Thật ra thì, hoa anh đào thuộc chi Prunus, có hơn bốn trăm loài khác nhau, nhưng ở Nhật thì phổ biến nhất là Somei Yoshino. À, các cậu có biết không, thời gian nở rộ của chúng chỉ kéo dài chừng một tuần, mà mỗi cánh hoa chỉ có năm cánh thôi đấy! Thú vị đúng không?''

Cô thao thao bất tuyệt, giọng đầy say mê. Shino-san thì đứng đơ ra, miệng mở mà chẳng biết phản ứng sao cho phải. Còn tôi chỉ kịp giơ tay, ngăn lại: ``Rồi, rồi, hiểu rồi, hoa đẹp lắm, Saki-san. Đủ thông tin để làm bài tiểu luận rồi đấy.''

Một tiếng cười khẽ vang lên phía sau. Yukihara Touka-san, cô bạn mà tôi nhớ là ``thần đồng'' của khối, đang cười nhẹ, như thể đã quá quen với tính cách mê hoa của cô ấy: ``Xin lỗi các cậu nhé. Cứ mỗi lần nói đến chủ đề hoa là cậu ấy quên mất cả thế giới xung quanh luôn\dots\ Mong các cậu thông cảm.''

\img{e3-8f}

``C-Chúng mình ổn mà\dots'' Suzu đáp lại, đổ mồ hôi hột, nụ cười gượng gạo.

Saki-san chẳng có vẻ gì là ngại, còn chỉ tay về phía ngoài sân: ``Đây này, các cậu nhìn đi!''

Ba chúng tôi đồng loạt hướng mắt ra ngoài. Dưới sân trường, hàng cây anh đào nở rộ, tán hoa dày đặc như mây. Gió thổi, hoa bay lên, tan vào ánh sáng như tuyết hồng.

``Oa\dots\ Đẹp thật đấy\dots'' Shino-san thì thầm, ánh mắt dõi theo những cánh hoa đang rung rinh trước làn gió.

Hai chị em tôi cũng không khỏi thừa nhận: đúng là một khung cảnh khiến lòng người nhẹ bẫng.

Cô bạn mê hoa nói tiếp, giọng chậm lại: ``Các cậu biết không? Ở thị trấn này, trước đây chẳng có nổi một cây anh đào nào đâu.''

Tôi bất giác thả chậm nhịp thở. Tôi cảm thấy như có bàn tay vô hình bóp nhẹ lấy trái tim mình: đôi mắt tôi mở to, đồng tử co lại khi hình ảnh những vòng lặp trước ùa về. Lúc đó, hai bên đường chỉ là gốc anh đào xám xịt, cành khô quắp như xương người giơ ra khỏi mộ; không một cánh hoa, chỉ có gió lạnh thổi qua khe nứt khiến tôi rùng mình.

``Ban đầu mình cũng không tin nổi đâu,'' Touka-san nói, như đồng tình với cô bạn mê hoa. ``Việc những cây hoa đào nở như vậy đã là một kỳ tích rồi đấy.''

Saki-san sau đó kể lại, tiếp lời: ``Ngày xưa vùng đất này bị nguyền rủa. Cỏ cây chẳng mọc nổi, con người cũng chẳng sống được.'' Lời cô ấy như gió thổi qua tai, nhưng trong tôi nó trở thành cơn bão. Tôi thấy chúng tôi đã từng đứng giữa phiên bản Aogami hoang tàn nhất: những tòa nhà bê tông nghiêng ngả, dây điện rối như mạng nhện giữa không trung, và trên mặt đường vỡ, những mảnh kính phản chiếu khuôn mặt tôi lúc bấy giờ: hoang mang và sợ hãi, nhất là sau khi chúng tôi lấy lại được ký ức của bản thân.

``Nhưng rồi, nhờ vào bốn con người cùng vị thần hộ mệnh của thị trấn này, với cả sức mạnh đến từ một thế lực bên ngoài nào đó, thị trấn đã được hồi sinh. Và hoa anh đào từ đó mới có thể nở trở lại.''

Tôi cảm thấy sống lưng mình khẽ lạnh. Bốn con người\dots\ sức mạnh ngoại lai\dots\ và vị thần hộ mệnh ư?

Shino-san quay sang tôi, ánh mắt cô lóe lên, sự kinh ngạc, xen lẫn nỗi xúc động khó tả. Còn Suzu thì siết chặt hai tay, không nói lời nào.

Những gì Saki-san vừa kể trùng khớp hoàn hảo với những gì đã xảy ra. Chính chúng tôi đã gieo hạt giống đầu tiên, khôi phục Aogami sau lần cuối cùng nó sụp đổ. Có lẽ nào\dots

``Hoặc ít nhất,'' Saki-san cười, vỗ nhẹ tay, ``đó chỉ là lời đồn thôi\~{}''

Cả ba chúng tôi thở phào. Nhưng trong khoảnh khắc ngắn ngủi đó, không ai nói gì thêm. Chúng tôi chỉ im lặng nhìn hoa rơi, như để cảm nhận lại điều mà có lẽ cả thị trấn này cũng đã quên mất: rằng đôi khi, huyền thoại bắt đầu từ chính những người bình thường.

``Cậu đừng có mà đi rêu rao mấy tin đồn linh tinh nữa đấy,'' Touka-san lên tiếng, giọng nửa trách nửa cười.

``Rồi, rồi, tớ biết rồi\~{},'' cô bạn mê hoa lè lưỡi, vẻ hối lỗi giả vờ, khiến tất cả chúng tôi bật cười. Sau đó, cả hai quay về lớp.

Tôi ngẩng đầu nhìn theo, rồi vô thức hướng mắt ra ngoài sân. Gió lại thổi, và một cánh hoa đáp xuống ngay cạnh tấm biển gỗ dưới gốc cây lớn nhất. Dòng chữ khắc nguệch ngoạc nhưng vẫn rõ ràng: ``Hoa anh đào Shinomiya.''

\img{e3-8g}

Cả ba chúng tôi sững lại.

``Shinomiya\dots?'' Shino-san khẽ đọc thành tiếng.

Ngay lập tức, một giọng nói lanh lảnh vang lên từ sau lưng, đầy vẻ tự mãn: ``Công lao của tôi đấy.''

Cả ba chúng tôi giật mình quay lại.

Sakura-san đang đứng đó, váy đồng phục khẽ bay trong gió, đôi mắt ánh lá mạ lấp lánh như cánh hoa rơi. Cô chống hông, vẻ mặt đắc thắng như vừa hoàn thành một kỳ tích: ``Không tệ chứ? Nhờ tôi mà chỗ này mới có mấy hàng anh đào đẹp thế này đấy.''

Tôi khẽ nhăn mày, nửa trêu nửa thật: ``Phải rồi, thần hộ mệnh của thị trấn mà. Cái gì cũng đến tay cậu cả. Với cả, đừng có nói mấy chuyện kiểu này trước mặt người khác chứ. Lộ ra là mệt đấy.''

Thần hộ mệnh nhún vai, vẫn giữ nụ cười tinh nghịch: ``Biết rồi\~{} Biết rồi mà\~{} Miễn là tôi không nói, ai mà biết được, đúng không?''

Shino-san thở dài, vừa cười vừa trách nhẹ: ``Từ khi nào mà cậu lại trở nên bất cần đời như vậy\dots''

Suzu thì chỉ khẽ lắc đầu: ``Nhưng mà\dots\ đúng là đẹp thật. Dù biết người làm là cậu, mình vẫn thấy khó tin đấy, Sakura-san.''

Cô nàng tóc hồng đỏ khẽ nghiêng đầu, ánh mắt dịu đi: ``Vì đó là điều duy nhất tôi có thể làm để `chuộc lỗi' cho vòng lặp trước thôi mà.''

Khoảnh khắc đó, gió lại thổi mạnh hơn, cánh hoa bay vòng quanh bốn đứa chúng tôi. Shino-san mỉm cười, nhỏ nhẹ: ``Thế này là đủ rồi. Chỉ cần thế này thôi.''

Sakura-san khẽ gật, rồi chậm rãi xoay người đi về phía lớp, bóng cô dần hòa vào những cánh hoa rơi, tan vào trong ánh nắng.

Tôi nhìn theo, thở dài: ``\dots Giữ bí mật, nhé.''

Em tôi gật đầu, Shino cũng khẽ đáp: ``Ừ. Mình nghĩ\dots\ ta chỉ nên giữ nó như là một `lời đồn' thôi.''

Sau đó, khi chuông báo hết giờ nghỉ vang lên, chúng tôi cùng nhau quay về lớp, mang theo trong lòng một cảm giác khó tả.

Chúng tôi biết rằng, dù Saki-san gọi đó là ``lời đồn'', nhưng năm chúng tôi đều hiểu rõ rằng chính chúng tôi mới là người viết nên truyền thuyết của thị trấn này. Và đôi lúc, sự thật về chúng nên được giữ kín.

\ct

\pov{Hikawa Kuon}

``Phù\dots\ Mệt bở hơi tai luôn.''

Tôi vừa mới giảng xong bài cho hai cô bạn ngồi ở bàn đầu - Kana-san và Kanade-san. Cả hai đều đang hí hoáy ghi chép lại những gì tôi vừa vẽ minh họa lên giấy nháp. Mà phải nói thật, dạy cho hai người có tính cách trái ngược nhau như thế cũng mệt ngang làm bài thi. Cô bạn Ác quỷ thì luôn chăm chú, ghi từng công thức, từng ví dụ, còn cô bạn Thiên sứ thì vừa nghe vừa gật gù, có khi còn ngáp dài rồi lại mỉm cười ngây ngô.

Khi họ rời đi, tôi nằm ườn xuống bàn, thở ra một hơi mệt nhọc. Mới có buổi đầu tiên thôi mà cảm giác như vừa trải qua cả tháng học liên tục vậy.

Từ góc phòng, tôi nghe loáng thoáng giọng Saki-san đang kể chuyện với nhóm bạn.

``\dots Cậu biết không, người ta đồn rằng thị trấn Aogami từng được tái sinh nhờ bốn con người và vị thần hộ mệnh của họ. Cùng với một sức mạnh đến từ thế giới khác nữa đó.''

Tôi nửa muốn phì cười, nửa muốn nghe tiếp. Câu chuyện nghe vừa hoang đường, vừa\dots\ có gì đó quen quen.

Touka-san ngay lập tức chen vào, giọng như đang trách móc nhẹ: ``Cậu thôi bịa mấy thứ vớ vẩn ấy đi. Lần trước cậu nói thị trấn bị ma ám, giờ lại thành được hồi sinh à?''

Saki-san cười xòa: ``Thì nghe người ta kể lại thế thôi\~{} Tin hay không là tùy cậu mà!''

Tôi chống cằm, nhìn ra cửa sổ. ``Bốn con người và vị thần hộ mệnh\dots'' câu nói đó khiến trong đầu tôi thoáng qua những hình ảnh mờ nhòe: một ngôi đền trắng xóa trong màn sương, lời cầu nguyện run rẩy của Shino-san, và ánh sáng rực chói nuốt trọn mọi thứ, ngay sau vòng lặp quá khứ đầu tiên.

``Thần linh, xin hãy cứu bọn con\dots''

Tôi khẽ thở dài.

``Nếu mà chỉ đơn thuần cầu nguyện những điều tốt đẹp như vậy thì nghe cực kỳ vô lý. Không làm mà đòi có ăn thì có mà cạp đất ăn à\dots'' tôi vô thức lẩm bẩm một mình.

``Tôi đồng ý với suy nghĩ đó đấy, Kuon-san ạ. Nhưng mà từ ngữ cậu dùng\dots\ hơi thiếu tế nhị đấy.''

Giọng nói vang lên từ phía sau. Tôi ngoảnh đầu lại, bắt gặp một ánh nhìn lấp lánh hai màu - vàng và đỏ. Một cô gái tóc nâu-đỏ buộc hai bím, dáng vẻ vừa sắc sảo vừa nhẹ nhàng.

\img{e3-8h}

``Akagane Mari-san\dots\ cậu là người của câu lạc bộ Mỹ thuật, đúng không nhỉ?''

Mari-san nghiêng đầu, nụ cười thoáng qua môi: ``Nếu mà chỉ cầu nguyện suông thì chưa đủ đâu. Thế giới này sẽ trở nên vô nghĩa nếu cậu chỉ thay đổi nó bằng cách mong ước nó thôi.''

``Đương nhiên rồi.'' Tôi bật cười nhẹ, mắt vẫn hướng ra ngoài cửa sổ. ``Nói mồm suông thì ai mà chẳng làm được. Ước là một chuyện, còn hành động hay không lại là chuyện khác.''

Mari-san gật đầu, ánh mắt thoáng pha chút đồng tình. ``Nếu cậu muốn thay đổi điều gì đó, thì phải hành động một cách thiết thực. Dù chỉ là một bước nhỏ, nó vẫn hơn gấp trăm lần so với việc đứng yên.''

Tôi cười nhẹ. Câu nói ấy\dots\ nghe có vẻ lý tưởng, nhưng lạ thay, tôi cảm thấy nó đúng thật. Trong lòng tôi bỗng nhớ lại Shino-san, người đã không chỉ cầu nguyện, mà còn đứng dậy, chiến đấu, và lựa chọn con đường của mình, tin tưởng vào ba chúng tôi, những người tới từ thế giới khác. Có lẽ nhờ vậy mà tất cả mới được cứu.

Tôi chọc nhẹ: ``Tôi không nghĩ một người yêu cái đẹp như cậu lại có thể nói những câu triết lý sâu sắc thế đấy.''

Mari-san khẽ mỉm cười, đáp gọn gàng mà có phần tinh nghịch: ``Nghệ thuật cũng là hành động, không phải sao? Vẽ một bức tranh đẹp đôi khi còn khó hơn cả thay đổi thế giới đấy.''

``Ừ\dots\ chắc vậy.'' Tôi gãi đầu, cười nhạt. ``Nhưng nói thế chứ\dots\ cậu tìm tôi có chuyện gì à?''

Cô nàng chớp mắt, nghiêng đầu ra vẻ hối lỗi: ``À\dots\ thật ra thì\dots\ tôi đang bí ở bài Toán phần đạo hàm ngược. Cậu có thể chỉ giúp không?''

Tôi lặng người trong ba giây, rồi thở dài, chống trán: ``Cuối cùng lại đi vào trọng tâm vấn đề rồi nhỉ\dots''

Cô ấy bật cười, cúi đầu xin lỗi, và tôi đành chịu thua. Thôi kệ, dù gì tôi còn năm phút nữa là hết giờ nghỉ. Thêm một người nữa chắc cũng chẳng chết ai đâu nhỉ?

\ct

Khi Mari-san quay về chỗ, tôi liếc nhìn quanh lớp. Ở bàn gần cửa sổ, Kana-san và Kanade-san đang bàn nhau về việc dựng kịch bản cho bản tin trường chiều nay.

\img{e3-8i}

``Hay mình làm vở kịch học đường kiểu `tình bạn xuyên thời gian' đi, nghe có vẻ thú vị đấy chứ!''  cô bạn Thiên sứ đề xuất, đôi mắt long lanh.

Cô bạn Ác ma chau mày: ``Lần trước bà viết kịch bản mà bọn mình phải quay lại ba lần đấy.''

``Thì lần này khác chứ\~{} Tôi đã học hỏi từ lỗi lầm rồi mà!''

``Câu đó bà nói tuần trước cho buổi nhập trường cũng y chang như vậy đó, Kanade.''

``Ờ\dots\ ờ thì\dots'' Kanade-san cười trừ, gãi má.

Bên kia, ở góc lớp, Mitsuki-san đang nghịch tóc Yuki-san, trong khi Hotaru-san ngồi quay bút trên tay, giọng như loa phường.

\img{e3-8j}

``Này này, tan học xong đi đâu nè? Tôi đăng ký đi hát karaoke nha\~{}'' cô bạn tóc nâu nói, ánh mắt sáng rực.

``Đi dạo trung tâm mua sắm đi, tôi thấy có tiệm kem mới mở đó.'' cô bạn tóc hồng chen vào.

``Cả hai đều sai rồi! Phải đi xem diễu hành lễ hội mùa hè chứ!'' cô bạn da ngăm nói chắc nịch, vung tay mạnh đến mức cây bút bay thẳng ra ngoài cửa sổ.

``\dots Ê, bay mất rồi kìa na,'' Mitsuki chỉ ra ngoài. Hotaru ngẩn người, rồi kêu toáng lên: ``\textbf{Á á á! Cây bút của tôi!!}''

Tôi phì cười. Ừ, đúng là ``thế giới thực'' thật rồi. Không còn những vòng lặp vô tận, không còn tiếng chuông trường báo tử, không còn những khuôn mặt đau đớn hay méo xệch. Chỉ còn lại âm thanh của tiếng cười, của ánh nắng rọi qua khung cửa sổ, và cảm giác bình yên đến lạ.

Tôi khẽ tự nói với mình: ``Có lẽ\dots\ đây chính là quả ngọt mà cho cuộc hành trình của chúng ta trong suốt thời gian qua.''

Tôi tựa đầu vào tay, mắt khẽ lim dim. Lâu lắm rồi, tôi mới thấy mọi thứ đơn giản như thế này. Một ngày học bình thường, tiếng phấn viết trên bảng, tiếng nói chuyện rì rầm, ánh sáng tràn qua cửa sổ, những thứ mà trước đây tôi từng xem là hiển nhiên, giờ lại trở nên quý giá đến lạ. Không chỉ bởi việc chúng tôi đã trải qua một chặng đường dài để tới được dây, mà do bản thân tôi, tất cả những thứ đó đã qua đi kể từ khi tôi bắt đầu bước chân vào giảng đường đại học và bắt đầu đi làm từ hơn hai năm trước, vậy nên, được sống lại khoảng thời gian đẹp nhất của đời học sinh dù gì cũng không phải quá tệ.

Tôi từng nghĩ ``tái sinh'' là một phép màu vĩ đại, thứ phải đổi bằng máu, nước mắt, hay hy sinh to lớn nào đó. Nhưng bây giờ, khi nhìn thấy Shino-san cười rạng rỡ với các bạn cùng trang lứa, thấy Kaigou-chan nói chuyện với những người bạn mới, hay nghe tiếng cười của Suzuki vang lên giữa lớp học, tôi mới hiểu: tái sinh đôi khi chỉ đơn giản là được sống tiếp, được cảm nhận niềm vui vụn vặt của hiện tại.

Không phải là thế giới đã thay đổi.

Mà là chính chúng tôi đã thay đổi.

Tôi mỉm cười, khẽ lẩm bẩm như lời kết cho dòng suy nghĩ đó. Có lẽ đây mới là điều mà Thần linh thật sự mong muốn.

\ct

\pov{Misora Shino}

Tôi cùng Sakura-san, Kaigou-san và Suzuki-san trở về lớp. Ánh sáng bên ngoài vừa đủ chan hòa, khiến cả hành lang sáng rực lên một màu ấm áp của buổi giữa xuân. Chúng tôi vừa đi vừa trò chuyện, về những người bạn vừa làm quen. Thật kỳ lạ, khuôn mặt thì quen thuộc đến mức tôi có thể gọi tên từng người, nhưng giọng nói, ánh mắt, cách cười của họ lạikhác hẳn. Có sức sống hơn, thật hơn, như thể những con người trong mảnh ký ức xưa kia đã được thổi hồn một lần nữa.

Trong khoảng mười phút làm quen ấy, tôi đã nghe đủ chuyện: từ việc Hanabi-san sợ trễ giờ, đến chuyện Inari-san làm rơi tập vở rồi cứ cúi đầu xin lỗi không ngớt, hay Saki-san và Touka-san nói về hoa anh đào như thể đó là điều đẹp nhất trần đời. Thật khó tin\dots\ những người từng chỉ tồn tại như những ``mảnh ký ức đã phai'' giờ đây lại đang sống thật sự, ngay trước mắt tôi.

Chúng tôi vừa về đến lớp thì cánh cửa trượt bật mở, một thứ âm thanh ấy vang lên rõ đến lạ. Kaoru-san bước vào trước, trên tay cầm hộp sữa, tay còn lại thì xách theo Miku-san, người vẫn đang lơ ngơ chưa hiểu chuyện gì đang xảy ra. Theo sau là Uzuki-san, đang cầm một chiếc máy ảnh kỹ thuật số sáng loáng, và Saya-san thì hí hửng vừa đi vừa cười.

\begin{figure}[H]
  \centering
  \begin{subfigure}[t]{0.45\textwidth}
    \centering
    \includegraphics[width=8cm]{../../../../Image/Book?/e3-8k1.png}
  \end{subfigure}
  \quad
  \begin{subfigure}[t]{0.45\textwidth}
    \centering
    \includegraphics[width=8cm]{../../../../Image/Book?/e3-8k2.png}
  \end{subfigure}
\end{figure}

``\textbf{Nào, cả lớp ơi!}'' cô bạn tóc ngắn hô vang, giọng cao vút, ``\textbf{Câu lạc bộ Nhiếp ảnh muốn có một kiểu ảnh cho lớp đây!}''

``Ể\~{}? Sao lại là lúc này chứ?'' cô bạm bờm Sói càu nhàu, hai chân vẫn khua trong không trung vì bị cô bạn tóc ngắn giữ chặt.

``Chụp thì cứ chụp đi, lo gì\~{}'' cô nàng tóc vàng ngân giọng, vẫn giữ nụ cười tươi rói. ``Rồi, rồi, tất cả mọi người vào vị trí nào!''

Cả lớp lập tức náo động. Ghế bàn xê dịch loạt xoạt, tiếng giày dép lạo xạo trên sàn gỗ. Tất cả đổ dồn về phía bục giảng, người thì ngồi hẳn xuống, người thì đứng chụm vào, người lại leo lên trên cùng để không bị khuất.

Kaoru-san chống tay hông, tuyên bố hóm hỉnh: ``Lớp mình hôm nay vắng sáu người nha\~{} Ayame, Shion, Aku, Koishin, Wata và Azuki, coi như chịu thiệt nha!''

Uzuki-san đứng bên dưới, đang chăm chú căn chỉnh máy ảnh.

``Ủa, sao cậu không lên chụp cùng mọi người đi?'' giọng Kuon-san từ hàng ghế sau vang lên.

Cô bạn tai thỏ giật mình, quay lại: ``Thì\dots\ tôi là người chụp mà, đúng không?''

``Dù gì cậu cũng là thành viên trong lớp 1-C mà,'' cậu nói, giọng đều đều nhưng ánh mắt lại pha chút vui đùa. ``Để tôi xử lý phần chụp ảnh này cho.''

Uzuki-san nheo mắt nghi ngờ: ``Cậu có biết chụp ảnh không vậy?''

Kuon-san nhún vai: ``Mẹ tôi từng bảo tôi có năng khiếu chụp ảnh. Thề, không điêu. Cứ tin ở tôi đi.''

Cô ấy ngần ngừ vài giây, rồi cuối cùng cũng chịu rời khỏi vị trí. Cô leo lên bục giảng, ngồi bên cạnh Yuka-san, người đang gật gù như sắp ngủ gật.

``Còn hai người nữa,'' cậu nói, quay sang phía hai chị em. ``Nếu muốn thì lên chung đi, tôi lo phần kỹ thuật cho.''

Suzuki quay sang Onee-sama của mình, chỉ thấy cô nhún vai, nở nụ cười: ``Lên thôi chứ nhỉ. Dù sao cũng là kỷ niệm.''

Hai chị em bước lên, đứng ở rìa ngoài cùng cạnh Kana và Kanade, trong khi cậu bắt đầu căn chỉnh hàng.

``Yae-san, cúi xuống chút\dots\ đúng rồi. Nanase-san, hơi nghiêng sang trái. Akane-san, đánh thức bạn cậu đi, chứ kẻo cô ấy ngủ thật đấy! Sakura-san, dịch sang trái hai bước-- không, không, sang trái của cậu cơ!''

Cả lớp cười ầm lên. Nhìn cách cậu căn chỉnh mọi người, trông chẳng khác gì một nhiếp ảnh gia chuyên nghiệp, vừa lạnh lùng vừa tỉ mỉ đến khó hiểu. Khi mọi người đã ổn định, cậu lùi ra sau, giơ tay lên cao.

``Rồi nhé! Cười tươi lên nào! Một\dots\ hai\dots\ \textbf{ba này!}''

*TÁCH!*

Đúng khoảnh khắc đó, một luồng gió xuân thổi vào qua khung cửa sổ mở, kéo theo hàng trăm cánh hoa anh đào bay lả tả, phủ kín khung hình. Một vài cánh đậu lên vai tôi, vài cánh khác xoáy thành đường cong trên tóc Sakura-san, rồi tản ra như những đốm sáng nhỏ. Tôi nghe tiếng vài người xuýt xoa, còn Kuon-san thì chỉ cười, bấm thêm một kiểu nữa.

Sau vài kiểu, cậu nói to: ``Giờ mọi người cởi áo khoác và áo len ngoài đi nhé, để đồng bộ hơn!''

``Ể!? Nhưng vẫn còn hơi lạnh đó nha\dots'' Miku-san lên tiếng, giọng đầy lo lắng.

Cậu đáp lại, tỉnh bơ: ``Giờ này thì đủ ấm rồi đấy. Với cả, mọi người đứng sát vào nhau thì ấm cúng hơn mà, đúng không?''

Một vài tiếng cười vang lên. Dù hơi ngần ngại, cuối cùng ai nấy cũng đồng ý. Kuon-san hướng dẫn thêm:

``Để tránh lẫn áo, mọi người cứ đặt áo khoác lên bàn của mình nhé. Chứ cỡ người mấy bạn giống nhau thì dễ nhầm lắm.''

Cả lớp cười rộ. Khi ai nấy đã chuẩn bị xong, người cầm máy ảnh lại căn c hỉnh, ngắm khung hình, rồi ra hiệu: ``\textbf{Rồi, chuẩn bị! Một kiểu cuối cùng nhé!}''

Máy ảnh kêu *TÁCH* thêm vài tiếng, trước khi Kaoru-san bỗng hét lên ngay khi cậu định bàn giao lại máy ảnh cho chủ câu lạc bộ: ``\textbf{Ê, Kuon-kun! Cậu cũng phải vào chụp chung chứ! Không công bằng nha!}''

Cả lớp đồng thanh phụ họa: ``\textbf{Đúng đó! Vào đi\~{}!}''

Kuon-san lập tức xua tay: ``Thôi, thôi, tôi mà đứng giữa rừng con gái thế này thì kỳ lắm!''

``Không được phép từ chối đâu, Kuon,'' Shiina-san chống nạnh, mắt sắc như dao. ``Cậu mà bỏ chạy, tôi bắt cậu chụp riêng một kiểu hình sau giờ học đấy.''

Touka-san đẩy nhẹ gọng kính vô hình, giọng vẫn điềm đạc: ``Thống kê cho thấy lớp có ảnh tập thể đầu năm thì năm học sẽ ít biến cố. Cậu không muốn chúng tôi gặp xui chỉ vì cậu chứ?''

Mitsuki-san khẽ nghiêng đầu, cười hiền: ``Kuon-kun, chỉ một tấm thôi mà. Mọi người đều muốn có kỷ niệm đẹp với cậu\dots\ được không zậy?''

Sakura-san huýt sáo, vẻ châm chọc: ``Ôi, chàng trai bí ẩn sợ camera kìa. Sợ lộ cái cằm đẹp trai chắc?''

Tôi nhún vai, giọng trung lập nhưng đủ để cậu nghe rõ: ``Kuon-san, cậu có thể giữ mũ trùm đầu nếu muốn. Đừng khiến họ phải chờ thêm.''

Tiếng reo hò bùng lên, bảy tám cánh tay kéo cậu vào giữa khung hình.

Cuối cùng, sau sức ép của cả lớp, Kuon-san cuối cùng cũng chịu bỏ cuộc. Cậu thở dài, giơ máy lên như lá chắn cuối cùng: ``Được rồi, được rồi, tôi chịu. Để tôi đặt máy lên tripod\dots''

Cậu chỉnh bộ dụng cụ dựng máy ảnh mà Uzuki-san mang theo, căn khung hình, rồi chạy nhanh về phía bục giảng, chen vào khoảng giữa của Kaigou-san và Suzuki-san.

``\textbf{Rồi nhé!}'' cậu hô to, giọng đầy bất lực. ``\textbf{Một, hai, ba-- Cười nào!}''

*TÁCH!*

Máy ảnh lóe sáng.

Trong khung hình, cả lớp 1-C nở nụ cười rạng rỡ nhất giữa biển hoa anh đào đang rơi.

\il

Mặt trời đã lên cao, ánh sáng chiếu qua khung cửa sổ, rọi xuống lớp học 1-C như những vệt kim tuyến mềm mại. Mọi người vẫn còn cười nói rộn ràng về mấy tấm hình vừa chụp: Kaoru-san thì khoe là mình ``đẹp nhất trong khung hình'', Miku-san thì xấu hổ trốn sau lưng Saya-san, còn Uzuki-san thì vừa xem lại ảnh vừa cằn nhằn sao có người lại nhắm mắt giữa chừng.

Tôi ngồi ở bàn mình, khẽ tựa má lên tay, nhìn cảnh tượng ấy mà trong lòng dâng lên một thứ cảm giác thật lạ: vừa bình yên, vừa ấm áp, như thể thời gian đang ngừng lại chỉ để tôi có thể ghi nhớ tất cả.

Mọi người\dots\ đều đang mỉm cười thật lòng.

Không còn nỗi sợ, không còn sự hoang mang hay vòng lặp bất tận nào nữa.

Chỉ còn lại hiện tại, và hơi ấm của những con người này.

Tôi cúi xuống, trong tay là bức ảnh mà Kuon-san vừa in ra từ chiếc máy ảnh kỹ thuật số của Uzuki-san.

Cả lớp 1-C lúc đó, tất cả đều có mặt. Và cả cậu ấy nữa, đứng ở giữa, giữa hai hàng bạn bè, ánh mắt sáng, nụ cười dịu.

Bức ảnh ấy\dots\ trông thật sống động. Không phải là bức ảnh lưu giữ ``hình ảnh'' của một khoảnh khắc, mà là ``linh hồn'' của cả khoảnh khắc ấy.

Tôi khẽ siết tấm ảnh trong tay.

Trong khoảnh khắc đó, tôi nhận ra rằng\dots

Nếu như trước đây, tôi chỉ biết cầu mong một điều ước mơ hồ rằng ``mọi người hãy hạnh phúc'', thì giờ đây, tôi muốn biến điều đó thành một lời nguyện chân thành, một hành động thực sự.

Tôi muốn cùng mọi người bước tiếp.

Tôi muốn ở bên cạnh, nhìn nụ cười này, nghe giọng nói này, và lưu giữ chúng\dots\ lâu thật lâu.

Tôi nhắm mắt lại, hít một hơi thật sâu.

\begin{center}
  ``Xin hãy để những nụ cười này tồn tại mãi, để chúng ta - dù ở bất cứ nơi nào -  vẫn có thể nhớ rằng, đã từng có một lớp học như thế này, đã từng có những ngày tháng tuổi học sinh bình dị nhưng rực rỡ như thế này.''
\end{center}

Một cơn gió nhẹ khẽ thổi qua cửa sổ, mang theo vài cánh hoa anh đào lơ lửng rơi xuống bàn tôi.

Tôi nhìn chúng, nở một nụ cười nhỏ, rồi đặt tấm ảnh vào trong ngăn bàn, nhẹ nhàng như thể cất giữ một phần của trái tim mình ở đó.

Tôi nghĩ\dots

Nếu như đây là ``kết thúc'', thì có lẽ, tôi sẽ không còn sợ những khởi đầu mới nữa.

Bởi vì từ giờ trở đi, tôi không còn đơn độc.

Chúng tôi sẽ cùng nhau sống, cùng nhau mỉm cười, và tiếp tục câu chuyện ấy\dots

\dots mãi mãi.

\img{group2}

\dots

\dots

\il

\pov{Hikawa Kuon}

Giờ nghỉ trưa đã tới.

Ánh nắng xiên qua những tán cây, rọi xuống sân trường, nhuộm vàng những bậc gạch cũ. Tiếng nói cười, tiếng giày nện trên nền gạch vang lên đều đều, hòa vào tiếng ve vừa chớm hè.

Tôi bước ra sau trường, men theo lối mòn nhỏ dẫn đến mép rừng. Gió ở đây mát hơn, trong hơn, và tĩnh lặng đến mức chỉ nghe được tiếng lá xào xạc va vào nhau.

``\eng{Game,}'' tôi lên tiếng sau một hồi im lặng.

Ông ta vẫn chưa nói gì từ sáng đến giờ, có lẽ cũng đang bận quan sát. ``\eng{Lỗ xoáy đâu rồi? Lẽ ra nó phải xuất hiện rồi chứ?}''

Một khoảng lặng kéo dài. Rồi giọng Game vang lên, khàn khàn và méo nhẹ qua tín hiệu đồng hồ: ``\eng{Lạ thật\dots\ Tôi không tìm thấy sự xuất hiện của lỗ xoáy nào cả.}''

Tôi khựng lại. ``\eng{\dots Kể cả khi chúng ta đã khôi phục thị trấn này rồi ư?}''

``\eng{\dots Tôi không nghĩ chúng ta đã khôi phục hoàn toàn đâu, Kuon-user.}''

Câu trả lời đó khiến tim tôi đập lỡ một nhịp.

``\eng{Ý ông là sao?}''

``\eng{Nhớ lời kể của Sakura-san về chuyện những đứa trẻ bị chôn sống để làm lễ hiến tế nhằm ngăn sạt lở thuở xưa chứ?}''

Tôi gật nhẹ, dẫu biết ông ta không nhìn thấy. ``\eng{Có\dots\ Sao?}''

Giọng ông ta trầm hẳn xuống: ``\eng{Tôi vừa phát hiện vài thực thể\dots\ có cùng đặc điểm với những linh hồn như vậy. Chúng đang lang thang trong khu rừng này. Một vài hồn khác thì quanh quẩn ở thị trấn.}''

Tôi nhíu mày, cố giữ bình tĩnh. ``\eng{Chắc chỉ là ảo giác thôi.}'' Tôi liếc về phía đám lá đang rung nhẹ, ``\eng{Có thể do không khí ô nhiễm, hoặc tàn dư của vòng lặp khiến tôi tưởng tượng ra mấy thứ ma quái.}''

Game-san cười khẽ, đầy vẻ châm biếm: ``\eng{Kuon-user, cậu vẫn thích tự an ủi mình như vậy. Nhưng lần này không phải do ô nhiễm, cũng chẳng phải ảo giác đâu. Là thật đấy.}''

Tôi siết nhẹ cổ tay, cảm thấy chiếc đồng hồ lạnh ngắt. ``\eng{Nhưng mà nếu đúng là linh hồn hiến tế, thì chúng đã chết từ mấy trăm năm trước rồi còn gì. Làm sao chúng vẫn còn ở đây?}''

``\eng{Vì thị trấn này chưa hoàn toàn được chữa lành,}'' ông ta đáp, giọng trở nên nghiêm nghị. ``\eng{Ánh sáng trắng chỉ tái tạo phần thân xác, nhưng những vết thương tinh thần, những oán niệm vẫn bám lấy không gian. Cậu đang nhìn thấy dấu vết của chúng, không phải ảo ảnh.}''

Tôi lẩm bẩm, nửa tin nửa ngờ: ``\eng{Tôi vẫn muốn tin rằng đó chỉ là do tôi tưởng tượng\dots\ Nhưng nếu ông nói đúng\dots\ vậy thì chúng ta đang sống chung với ma quái thật sự.}''

``\eng{Chính xác. Và cậu sẽ phải đối mặt với họ, sớm thôi.}''

Một luồng khí lạnh chạy dọc sống lưng tôi. ``\eng{Không thể nào\dots\ Ý ông là\dots\ những hồn ma đó vẫn còn tồn tại ở đây sao?}''

``\eng{Phải. Nhưng bình tĩnh nào, Kuon-user,}'' ông ta trấn an, giọng đều đều. ``\eng{Hiện tại chúng chưa gây nguy hiểm gì. Có thể chúng vẫn đang chờ đợi điều gì đó. Dù sao, cứ đề phòng thì vẫn hơn.}''

Tôi ngẩng đầu, nhìn qua rặng cây. Xa xa, trong khoảng trống giữa hai thân gỗ, có thứ gì đó như làn sương nhạt đang lơ lửng, tản ra rồi tan biến.

Không khí vẫn trong trẻo, nhưng đột nhiên, tôi cảm thấy thị trấn này, dù đã rực rỡ, tươi sáng, vẫn còn ẩn chứa điều gì chưa được khôi phục hoàn toàn.

``\eng{Có lẽ,}'' Game-san nói tiếp, ``\eng{giải quyết chuyện của những linh hồn đó sẽ là chìa khóa để thoát khỏi nơi này. Nhưng trước mắt, cậu cứ sống như một học sinh bình thường đi đã. Có vài người trong lớp có thể giúp cậu đấy.}''

Tôi chỉ biết thở dài, lẩm bẩm với bản thân: ``\vie{Có lẽ chúng ta sẽ bị kẹt ở đây lâu hơn nữa rồi\dots}''

Một thoáng im lặng. Gió thổi, mang theo hương cỏ ẩm và mùi phấn viết bảng còn vương từ dãy lớp học. Tôi ngước nhìn lên trời, ánh nắng xiên qua kẽ lá, nhòe trong mắt.

Không hiểu sao, trong khoảnh khắc đó, tôi lại nghĩ đến các cô gái ở thế giới gốc: A-san, Miko-chi, Sui-chan, Mio-chan, và cả\dots\ Robo-PON nữa. Không biết họ sẽ phản ứng ra sao khi tôi biến mất một cách đột ngột như thế đây.

Tôi khẽ mỉm cười.

``Thôi thì\dots\ sống lại đời học sinh một lần chắc cũng không tệ nhỉ?''

Tôi quay người lại, bước về phía sân trường. Xa xa, vọng lại tiếng gọi nhau í ới của mấy bạn cùng lớp, tiếng cười hòa lẫn tiếng gió.

Bước chân tôi nhẹ dần, hòa vào âm thanh ấy, như thể mọi thứ đang bắt đầu lại một lần nữa.

\end{document}